\quellseite{219}
\part[II. The Subject of Religious Instruction]{II. The Subject of Religious Instruction: The Pupil and His Inner Life}
\quellseite{221}
\chapter[§ 21. The Inner Life of the Pupil]{§ 21. The Inner Life of the Pupil in Its General Aspect}

J. F. Herbart, \emph{Lehrbuch der Psychologie}, 1816.\footnote{Arranged according to schools.} V. Lazarus, \emph{Das Leben der Seele in Monographien ueber s. Erscheinungen u. Gesetze}, 1856, ³1883. L. Struempell, \emph{Grundriss der Psychologie}, 1884. W. F. Volkmann, \emph{Lehrbuch der Psychologie}, 1876, ⁴1894. P. Natorp, \emph{Einleitung in die Psychologie}, 1888. J. Remhke, \emph{Lehrbuch d. allgem. Psychologie}, 1884, ²1905. H. Spencer, \emph{Principles of Psychology}, 1870. G. J. Romanes, \emph{Mental Evolution in Animal}, 1883. G. J. Romanes, \emph{Mental Evolution in Man}, 1888. Theod. Ribot, \emph{L’heredite psychologique}, 1873, ⁶1902; in German: \emph{Die Erblichkeit}, 1879. Th. Ribot, \emph{Les maladies de la memoire}, 1881, ⁴1898; in German: \emph{Das Gedaechtnis u. s. Stoerungen}, 1882. W. James, \emph{The \emph{Principles of Psychology}}, 2 vols., 1890. W. James, \emph{Psychologie, Shorter Course}, 1892. H. Muensterberg, \emph{Psychology and Life}, 1899. H. Muensterberg, \emph{Grundlinien der Psychologie}, 1900. E. Mach, \emph{Die Analyse d. Empfindungen u. d. Verhaeltnis d. Physischen zum psychischen}, ⁴1900. Th. Ziehen, \emph{Leitfaden d. Physiologischen Psychologie}, ³1908. F. Jodl, Lehrbuch d. Psychologie, 1909. W. Wundt, \emph{Vorlesungen ueber die Menschen- und Tierseele}, 1863, ⁴1906; English by J. E. Creighton und E. B. Titchener, 1894. W. Wundt, \emph{Grundzuege d. physiologischen Psychologie}, 1873, 1908. W. Wundt, \emph{Grundriss d. Psychologie}, 1896, ¹⁰1911. O. Kuelpe, \emph{Grundriss d. Psychologie}, 1893; English by Titchener, 1895. N. Ach, \emph{Die Willenstaetigkeit u. d. Denken}, 1905. A. Messer, \emph{Empfinden u. Denken}, 1908. E. Meumann, \emph{Vorlesungen z. Einfuehrung in die experimentelle Paedagogik}, 1907, ²1911. F. Meumann, \emph{Intelligenz u. Wille}, 1908. H. Ebbinghaus, Grundzuege d. Psychologie, 1897 ff., ³1911; English by M. Meyer. H. Ebbinghaus, \emph{Abriss der Psychologie}, ¹⁹10. E. B. Titchener, \emph{Textbook of Psychology}, 1910, 1912. E. B. Titchener, \emph{Psychology of Feeling and Attention}, 1909. E. B. Titchener, \emph{A Beginner’s Psychology}, 1915. James M. Baldwin, \emph{Elements of Psychology}, 1893. J. M. Baldwin, \emph{Development and Evolution}, 1902. K. O. Beetz, \emph{Einfuehrung i. d. moderne Psychologie}, 1907 f. Th. Elsenhans, Lehrbuch d. Psychologie, 1912. J. M. Baldwin, \emph{Handbook of Psychology}, 1889. 1891. \emph{Zeitschrift fuer Psychologie}. \emph{The American Journal of Psychology}. \emph{The Psychological}
\quellseite{222}
\emph{Review}.—J. F. Herbart, \emph{Allgemein. Paedagogik a. d. Zweck der Erziehung abgeleitet}, 1806. J. F. Herbart, \emph{Umriss paedag. Vorlesungen}, 1835; English by H. M. and E. Felkin, 1895. Compare John Adams, \emph{Herbartian Psychology applied to education}. T. Ziller, \emph{Grundlage z. Lehre v. erziehenden Unterricht}. 1865. O. Willmann, \emph{Didaktik als Bildungslehre}, ²1894 f. Fr. Froebel, \emph{The Education of Man}, translated by Hailmann, 1887. M. Jahn, \emph{Psychologie als Grundwissenschaft d. Paedagogik}, 1895, ⁶1911. P. Barth, \emph{Die Elemente d. Erziehungs- und Unterrichtslehre auf Grund d. Psychologie d. Gegenwart}, 1902, ²1908. W. Lay, \emph{Experimentelle Didaktik}, ³1910. G. Schumann u. G. Voigt, \emph{Lehrbuch d. Paedagogik}, vol. ii: Psychologie, ¹¹1901. Ostermann u. Wegener, \emph{Lehrbuch d. Paedagogik}, vol. i: Psychologie, ⁴1910. K. Heilmann, \emph{Handbuch d. Paedagogik}, vol. i: Psychologie u. Logik, ¹²1908. L. Hohmann, \emph{Paedagogische Psychologie}, 1906. L. Habrich (Roman Catholic), \emph{Paedag. Psychologie}, ⁴1911. K. Lange, \emph{Ueber Apperception, e. psychologisch-paedagogische Monographie}, ¹²1912; English by the Herbart Club and edited by Ch. de Garmo, 1894. James Sully, \emph{Outlines of psychology, with special reference to the theory of education}, 1894. P. Radestock, \emph{Habit and Its Importance in Education}, English by Caspari, 1887. D. Kay, \emph{Memory. What It Is and How to Improve It}, 1889. Th. Ribot, \emph{Psychology of Attention}, 1890. Jos. Baldwin, \emph{Psychology Applied to the Art of Teaching}, 1892. G. Compayré, \emph{Psychology Applied to Education}, 1894. W. T. Harris, \emph{Psychologic Foundations of Education}, 1897. Froebel’s \emph{Educational Laws for All Teachers}, 1897. W. James, \emph{Talks to Teachers on Psychology}, 1898. H. H. Horne, \emph{Psychological Principles of Education}. E. D. Starbuck, \emph{Psychology of Religion}, 1899. W. James, \emph{The Varieties of Religious Experience}, 1902. G. A. Coe, \emph{The \emph{Psychology of Religion}}, 1916.

The child or the pupil is the subject of religious instruction as imparted by the Church. Her object being his instruction and education, he must be accurately understood and the peculiarities of his inner life must constantly be kept in mind if the educational and training efforts are to be crowned with success.

The pupil consists of two essential parts—\emph{body and soul}, the former being his material, and the latter his psychical constituent. Materialism, indeed, denies the reality of the soul, explaining all psychic phenomena as mere products
\quellseite{223}
of physical or cerebral processes. However, certain facts of experience — such as the continuity of self-consciousness in the face of incessant organic changes, also in the brain; the unity of consciousness; the impossibility for a movement of material atoms to produce anything but another physical movement; the strife between the soul and the body and the dominion over the latter by the soul—and the Scriptures postulate behind the movements of the brain the existence of a power, invisible and independent, which permeates and determines the body and is essentially different from the brain and from all matter. It is this power which we call soul. In the pupil, then, two worlds essentially different from each other are welded together in wondrous union.

Between the body and the soul (the physical and psychical) there exists a vital reciprocity or interaction. Grief and mental anguish act upon the lacrymal glands; wrath, upon the liver; fear, upon the intestines; shame flushes the face with crimson and terror blanches it white. Conversely, a sojourn out-doors invigorates us mentally; lack of blood in the brain causes unconsciousness; and defective sense-organs inhibit the orientation of the soul to the world without. This interaction between body and soul is effected by the nerves. We distinguish two classes of nerves, the centripetal (centrum petere), and the centrifugal (centrum fugere) nerves. The former nerves, also called sensory or afferent, take note of stimuli\footnote{The stimulus may proceed from some object of the external world which acts upon the sense-organs, or from the body itself or one of its members.} and pass them on to the soul; the latter class of nerves, also called motor or efferent, forward the message of the soul to the various parts of the body. With lightning rapidity the excitation of a sense-organ is reported, by the sensory nerves, from the periphery to the brain; and upon receiving the message, the soul quickly executes its decision or choice through the motor nerves.—By a severance
\quellseite{224}
of the respective \emph{nerves} both sensation and motion may be totally inhibited, while they may be seriously disturbed and reduced in strength through congenital or acquired weakness or disease as is seen in the case of imbecility and neurasthenia.

Inasmuch as the soundness and the regularity of these interaction processes is of paramount importance, and since the pupil who is somewhat subnormal in mental efficiency requires special treatment, the teacher of religion must pay attention also to the physical welfare of his pupils, seeing to it that they receive pure air, adequate and correct food, proper clothing, and a due amount of rest and exercise. He must likewise insist on a proper posture in order to protect the sense-organs—the gates of the soul—, especially the eye and the ear, from injury. He must furthermore carefully study the pupils as to their nervous state. He must support the weakness of the feebleminded, who are, at least generally, quite capable of receiving sensory stimuli, but lack the faculty of prolonged concentration of attention on any one object and therefore acquire but a meagre amount of information. He must patiently bear with neurotic and neurasthenic children, never overtax their mental capacity and never fatigue them; he should arrange for frequent changes in the nature of their work, and apply some of the well known and sufficiently attested remedies. Above all, he must know the inner life of his pupils and govern his approach according to its peculiarities if he would really find the way to their inmost heart, take hold of their whole personality, and effect permanent results.

The soul is an indivisible unity. But it manifests itself in three different ways; the indivisible soul knows, it feels, and it wills. Accordingly we speak of the intellectual life, the emotional life, and the volitional life of the soul, all of which constantly interact and all of which must be stimulated if the whole personality is to be influenced.
\quellseite{225}
\section{1. The Intellectual Life}

The intellectual life of the soul is reflected in three stages, sensation, concept, and thinking. The first of these stages, in turn, has three aspects or steps, viz., sensation, perception, and intuition.

a. We experience a sensation when a stimulus proceeds from some external object, is received by the sensory nerves, and flashed to the brain. For instance, the light waves emanating from the sun touch the eye; the optic nerves terminating in the retina are stimulated thereby and transmit, as telegraph wires, the stimulus to the brain; thus the brain is stimulated by the nervous impulse, and as a result, the soul which must be carefully distinguished from the brain, is, in a manner unknown to us, set in motion, that is, it experiences a sensation. In other words, when a sensation is produced, three processes take place: a physical process, the external object meeting the eye and stimulating it; a physiological process, the stimulus being conveyed by the nerves to the brain where it makes a certain impression; and a psychical process, the soul receiving a sensation. Instead of from some external object, the stimulus may proceed from some state of the body, as hunger or fatigue; sensations thus arising are called bodily, or organic, sensations in distinction from the sense-sensations, which are brought about by the organs of sense. Sensation, whether external or bodily, is the first and the simplest activity of the intellect. Whatever may have taken place in the soul, there is no conscious soul-life until sensations have been registered. Corresponding to the several sense-organs the following classes of sensations are enumerated: light or visual sensations, sound or auditory sensations, these two being the main factors in the enrichment of our mental life; furthermore, smell or olfactory sensations, taste or gustatory sensations, and touch or skin or cutaneous sensations, the last
\quellseite{226}
three, like the bodily sensations, serving largely the interests of the animal life. All sensations, in turn, differ according to quality (as salt, sour, sweet, and bitter qualities of taste), intensity (a whisper barely crosses the threshold of stimulus, a cry may rise to the extreme height of intensity), and feeling tone (they may engender the concomitant phenomena of pleasure or displeasure).

In sensation, the soul is simply made aware of a certain stimulus. When, however, the soul isolates this sensation from others which may occur simultaneously, and traces it back outwardly and locates the object that caused it (whether this object be in the body or outside of it), that is to say, when the soul projects the sensation outwardly, then it forms a \emph{perception}, it has perceived the object from which proceeded the stimulus that caused the sensation. This paves the way for a further advance of the soul life. Now the soul can examine the object in question and employ all its senses in the exploration and investigation of the object in whole and in part (the senses of sight, touch, and smell are brought to bear upon yonder rose), and thus it gains a clear and complete image of the object. The image which is thus formed in the soul we call \emph{intuition}. The more completely the several characteristics of an object have been perceived, the more distinct will be the image formed by the soul; and the more plainly the object's distinctive qualities are noted, the clearer its image will be. Except upon the basis of distinct and clear intuitions there can be no further development of the mental life, no reliable knowledge, no correct judgment. For this reason it is of infinite importance that the teacher of religion should lead his pupils to form distinct and clear intuitions of things religious. Following the example of Luther, Ratichius, Comenius, and Pestalozzi, he must apply the principle of miniature painting and place before his pupils most vividly the outstanding exponents of the religious life
\quellseite{227}
so that they can not but grasp firmly what it means to be a Christian and obtain definite and clear-cut images of faith, love, fidelity, and the like. And as soon as he observes in the course of further instruction that the traces of such images are effaced in the mind of the pupil, and that he carries about lifeless concepts instead of living images of religious realities, he will at once return to the task of providing vivid intuitions; yes, he will not wait until this calamity has occurred but try to forestall it; at any rate, he will restore the vividness of intuition as soon as he discovers that it is becoming less definite and clear.

b. A higher step of mental life is reached through what is called concept.\footnote{The terms “concept” and “conception” are often used interchangeably. Here, however, we distinguish them from each other. We use the term “concept” in its original sense: what by outward means is brought into existence inside of a body; what remains in the inside when the outside factor is removed; the image remaining in the soul after the object producing it is taken away. The faculty of the soul to form and to retain such images in the mind is “imagination” understood in its original sense: the faculty whereby the “soul beholdeth the likeness of things that be absent.” By “conception,” however, we mean the first stage of the thinking process in the narrower sense, that is, that constructive act of the soul whereby two or more attributes are grasped into a unity of thought, and also the product of such act (cf. p. 215). “Conception,” as we use it here, then, is not possible without logical thinking, while “concept” precedes all logical thinking.} In forming intuitions it is required that the soul faces an object and perceives it by means of the sense-organs. Were there no further activity of the soul, it would resemble the mirror reflecting an image only as long as it is confronted by an object; in that case there could be no recollection of objects previously perceived. But the soul has the wonderful power of retaining the image of an object even after the object itself has been removed. The soul, as it were, places the object inwardly in front of itself, thus producing what is called a concept. This concept, indeed, lacks the vividness of the original sensation (the
\quellseite{228}
concept of tooth-ache does not cause pain) because the object from which the stimulus originally proceeded has been removed; nevertheless, the image will be as complete or defective, clear or confused, distinct or indistinct, as the preceding intuition was—another fact suggestive of the imperative necessity of the teacher’s providing complete, clear, and distinct intuitions. It is only through this faculty of forming concepts that a coherent mental life becomes possible for man, a unified thought life which does not grasp merely the fleeting moment but the past and the future as well, a compact continuity, a history of his life. Now there arises and takes form an inner world which constitutes the most characteristic feature of human soul life.

Even now, it is true, a coherent unity of mental life would be out of the question were it not for the faculty of the soul to put the concepts in motion, to associate them with one another and to recall such as have become obscure. But the soul possesses also this power; and accordingly we speak of the motion, the association, and the reproduction of concepts. A concept recently formed lives in the soul clear and strong; but when amid the perpetual changes of life ever new concepts are formed, then those previously formed do not retain their pristine clearness and strength, they wane and pale as new ones enter; or to use a figure, they are crowded back and sink beneath the threshold of \emph{consciousness} (\emph{consciousness} being the sum total of clear and strong concepts). In this way, the world of concepts is kept in constant motion, and we must distinguish between that of which the soul is conscious at any time and that of which it is \emph{unconscious}. For the soul can not at any moment retain all its concepts with an equal degree of strength and clearness, i.e., be conscious of them. Such power would far transcend the divinely appointed limits of the soul. As Kant says, “Our soul is like the vast treasure vault illuminated
\quellseite{229}
but dimly by a little lamp whose rays fall only on a small number of objects.” This “narrowness of consciousness” is another characteristic mark of the life of the human soul. In the state of sleep or in the swoon every trace of consciousness disappears. Yet, the soul loses none of the concepts ever formed in it; they persist in it, or rather, it retains them, even if beneath the threshold of consciousness. Nor does the soul allow new concepts to remain aloof from those that have been previously formed; it rather strives to connect or \emph{associate} the new and the old concepts. If the new concept has the same contents as the old one, both fuse into one which thereby obtains greater clearness. If the new concept is opposed to the old, there also ensues a fusion, but only by way of assimilation, the opposite elements being neutralized in the process. If, finally, the new concept is dissimilar or disparate to the old, the soul can not fuse them into one and must keep them separate, but it may gather the simple concepts (round, yellow, fragrant) into a total concept (rose). Or, the new concept is associated with previous ones at least in regard to space and time, that is mechanically, or logically. For the soul tends to link together and to connect closely all its concepts, and the more intimate such connection is, the more easily and tenaciously it retains them.

In the same manner also the \emph{reproduction} of concepts is assured, this being the next step in the development of mental life. The concepts may have sunk beneath the threshold of consciousness, yet the soul has the faculty of causing them, as it were, to rise again and can become conscious of them once more. Such \emph{reproduction}, however, in most cases does not occur immediately, but mediately. We speak of immediate \emph{reproduction} when a concept that has disappeared from consciousness, rises spontaneously into consciousness without the stimulation or cooperation of any conscious
\quellseite{230}
 concept; this takes place only if the concept when originally formed deeply aroused our interest and strongly affected our emotions. Mediate reproduction we term the process of recalling a concept to consciousness by means of another, conscious concept. The following laws govern such reproduction: (1) If the soul simultaneously formed a number of concepts, it recalls by means of the one concept which has remained in consciousness all those of which it has become unconscious (the dog whom I have always seen following a certain man recalls his image to my mind after he is dead). (2) Concepts which have successively entered consciousness recall one another in their original order (the initial measures of a tune recall the rest). (3) Similar concepts recall one another (the judgment visited upon Sodom and Gomorrah recalls that of the deluge). (4) Contrary concepts recall one another (the happy features of his home life were recalled by the prodigal son when he meditated upon the opposite, unhappy ones of his present condition).

Also these characteristics of the pupil’s mental life suggest to the teacher of religion a number of hints and duties. The “narrowness of consciousness” will prevent the teacher from offering in rapid succession a vast number of manifold things, from assailing the soul of his pupil with new objects while it is still engrossed with those previously offered. The power of the soul to retain and to recall all concepts once formed, serves to impress upon him the need of keeping as far from the pupil as possible everything sinful and vile, of supplying his soul with clear concepts of moral and religious nature, and of setting in word and deed an example of a true Christian life. The tendency of the soul to associate the new concepts with the old ones imposes upon him the duty of connecting each new cognition with the existing sum of knowledge in order that the pupil’s knowledge will become an organic unity and as such influence the soul and
\quellseite{231}
in order that later the one concept may lift the others into consciousness and lend support to them. And in view of the fact that the reproduction of former concepts is facilitated and assured by the interest that had been aroused when they were first presented to the pupil, and by the frequency of repetition (according to the law of fusion), he will endeavor, by every possible means, to arouse the pupil’s interest in religious materials, and he will never tire of careful and frequent reviews conducted at the right time (“If review becomes necessary, it is too late.”) and in the right manner, stressing especially the essentials and paying due attention to drill work. An understanding of the principles here involved will lead the teacher to appreciate the significance of religious habits and customs and of their cultivation from early youth.

The soul’s power of unaltered reproduction is called \emph{memory}. This can and should be strengthened by appropriate exercise; but it should not be overtaxed and weakened. When the soul purposely appropriates several concepts in definite order, it “memorizes.” Such memorizing is logical or judicious when the individual concepts to be learned are viewed from the standpoint of order and inner connection; logical memorizing facilitates appropriation, retention, and recall, and in general enhances one’s mental powers. If, however, the purpose is unaltered, verbatim, reproduction and absolute mastery of the material, mechanical \emph{memory} must be combined with the former. Occasionally a concept not germane to the material may be utilized to good advantage in the learning and recalling process; e.g., in order better to remember the beginning of the several stanzas of a poem, a sentence may be formed of the initial words or letters. This is what we call ingenious, or mnemonic, memorizing. Also in religious instruction the teacher is occupied with the appropriation of facts and truths, of verses and stanzas. Now
\quellseite{232}
he will certainly exercise the pupils especially in logical memorizing, but he will not overlook the value of mechanical memorizing in connection with the former. By all means, however, he will guard against the error of having his children memorize anything upon which no light has been shed as to contents and connection; otherwise memorizing will become a burden, and the work of memorizing which might have been of value to the pupil, brings him no gain—no matter how significant the memory material may be in itself and prove to be in the future.

The soul is not compelled to reproduce the acquired concepts in unaltered state; it has the power of changing them and combining them into totally new forms. It disposes of them with the same freedom with which the type setter rearranges his types in combinations ever new. It may eliminate and abstract, supplement, combine. This faculty of the soul is called \emph{phantasy} or imagination. The soul forms abstractions when it eliminates whatever is incidental and retains what is essential, when it conceives as independent of their object attributes and activities which adhere to certain objects and ordinarily cannot be thought of as existing by themselves.\footnote{The attribute round and the activity of motion I cannot conceive with my senses as real and independent of any object. They rather coalesce (\emph{concrescere}) with the object, in this case the globe, a fact which has given rise to the name assigned to such objects, viz., \emph{concretum}. However, the soul is able to detach its attention from the object as such (\emph{abstrahere}); to note only the roundness and rolling movements of the globe, thus producing the concept roundness.} Again, the soul is able to add to the abstract and general the concrete and the specific, and thereby to supplement, to illustrate, and to vivify the abstract. Finally, it can take existing concepts and by variously combining them create new forms. By virtue of \emph{phantasy}, then, the soul becomes an artist creating a new world, a world of its own on basis of existing things and through powers inherent
\quellseite{233}
in itself. Phantasy makes the inner life rich and varied; it familiarizes the soul with the past and it permits glimpses into the future; it renders the soul sympathetic with the weal and the woe of others, participant in the activities of their life, willing to help. For, as the \emph{phantasy} stirs the emotions through vivid images, so it also arouses the will to action. There is every reason, then, that the teacher of religion should quicken and nurture the \emph{phantasy} of his pupils by supplying many details and by cultivating clear intuitions. Failing to do so, he would leave unused the key to the emotions and to the will of his pupil; he would neglect the opportunity of leading him to thrilling joy, to sincere grief, to noble aspirations, to ardent enthusiasm for all that is good, true, and beautiful; he would certainly have failed to learn from Holy Scripture which, fairly brimming with pedagogic wisdom, never wearies, e.g., to paint the glory of eternal life in ever new and ever changing images.—It is hardly necessary to call the teacher’s attention to the fact that, despite many points of similarity, healthy \emph{phantasy} is not identical with the uncontrolled imagination. The former requires cultivation, the latter repression; for, delusions, illusions, and hallucinations are “fantastic” mental products having their origin in a departure from the normal processes of perception and intuition.

Owing to their preservation and reproduction by memory, however, the earlier concepts do not only constitute the material for the creative \emph{phantasy}, but they provide also the means to interpret and to understand new objects presenting themselves to the soul. When the child who has already become acquainted with some birds, for the first time sees a crane, he does not only receive an image transmitted through his sense-organs, but he is also at once more or less distinctly reminded of the similar birds which he has previously observed; spontaneously he compares with them this new
\quellseite{234}
animal, perceives the points of evident resemblance and of dissimilarity, and thus perceives the crane as a bird. We use the term \emph{apperception} (\emph{ad-percipere}) for this power of the soul to interpret new concepts by means of old ones, to understand them, and correctly to associate them with concepts it already possesses. Apperception is of the utmost importance. Every perception of a definite object in the external world; every sensory cognition and recognition of things, persons, or phenomena; all comprehension of words heard or read; the interpretation of any figure of speech—all takes place through this illuminating association of the new concepts with the old ones, that is to say, through \emph{apperception}. The soul, as it were, leans upon the old concepts, it employs them as the basis from whence it may secure mastery over the new material and properly classify it. Old concepts in the soul serve as the key to the new concepts which may be appropriated either for its own enrichment or for the clarification and correction of its possessions, for the new when once it is cognized, in turn sheds light upon the old so that learning is not only an acquisition, but frequently also a transformation, of knowledge.

A principle of great didactic significance, also for religious instruction, results from the fact of \emph{apperception}, viz., the principle: “From the old to the new, from the known to the unknown, from the near to the remote!” Apperception places upon the teacher the sacred duty to acquaint himself with the concept mass at the command of his pupils. He will obtain this information by his fellowship with them, by observing them at play, by watching the literature which they read, by visiting their homes, perhaps also by requiring written work now and then, and by paying incessant and careful attention to the questions and the answers which they offer during the process of instruction. When he has thus gained an insight into the concept world of his pupils,
\quellseite{235}
he will carefully select those elements which may serve as points of contact for the new things to be offered, and let them shed the proper light upon the new material. Thus he will assist the pupils in progressing step by step without gaps or leaps to thorough understanding, coordination and mastery of the new material, or possibly to the clarification and correction of the old. “For what one has no related thoughts, no points of contact, no apperceiving material, for that one has no eye, no receptivity, no understanding.” If the teacher of religion does not come down in his language to the vocabulary of his pupils, or if he speaks of objects altogether beyond their mental horizon, then the children may have some auditory sensations, but their minds do not understand or they misunderstand. The reason is that in their concept realm there is no concept that might serve as a clew to the mysterious words or images employed by the teacher. As an inevitable result, the pupils lose \emph{attention} and \emph{interest}, and abandon all hope of ever mastering the new material. Conversely, the teacher who interprets the new in the light of the old, arouses his pupils’ \emph{attention}, i.e., he focuses their mental activity on the one point from which he would lead them to the comprehension of the new, and now inaugurates a process of observation and comparison between old and new, of joint exploration and investigation. Every discovery, i.e., every step of progress upon the way to complete mastery, inspires them with new joy and courage for further effort, he has aroused a vital \emph{interest} in the task: his pupils gladly occupy themselves with it, love it, are engaged in it with all their soul, and never desist until they have solved the problem and mastery has crowned their efforts.

c. Thinking forms the third stage of development in the intellectual life of the soul. In its elementary form, denoting all conscious intellectual activity, thinking is found
\quellseite{236}
already in sensation and, even more, in concept. Except for the thinking activity the soul, having received a stimulus, could not possibly trace it back to the object causing it nor gain a clear intuition of it; nor would it be able to recognize concepts or associate and combine them; nor would it be able, by the use of phantasy, to dispose freely of the concepts in its possession; nor could there be any apperception of the new concepts by means of the old ones. From this more or less elementary type of thinking, however, we must distinguish the \emph{logical type} which marks a distinct advance over the lower forms. The soul is able not only to perceive objects, to understand relations of time and space, and to observe other incidental relations, but also to cognize the inner, or essential, relations, that is, it can determine logical connections. It examines the objects and concepts that come within its sphere of observation as to their inward relation (cause and effect; \emph{reason} and result; means and purpose; genus and species; object and attribute, etc.); it associates what belongs together, and separates what does not; thus it brings order and clearness into its store of mental treasure and is enabled to infer new facts from those already appropriated and clearly understood. This activity of the soul we call logical thinking, and the faculty exercising it, \emph{reason}. Thus, when, in physics, one phenomenon is recognized as caused by a natural force; when, in history, one event is shown to be the effect of previous ones; when, in language, the nature of an element of the sentence is demonstrated at the hand of a large number of examples or a poem is analyzed in regard to connection or progress; when the motives back of a certain action are investigated — always the object is to ascertain inner or logical relations between the concepts in question, that is, logical thinking.

If it is true that actual, complete control of the world of concepts is not achieved except by logical thinking, how
\quellseite{237}
could the teacher of religion possibly fail to cultivate it? He will not be satisfied to treat unconnectedly the individual stories of the Bible, or to enumerate mechanically the several features of any one story; he will rather induce the children to inquire after the inner connection and the impelling motives, after ground and consequences or cause and effect. He will take pains to set forth the underlying principles of divine and human action which are still in force today and clamor for observance, no matter how widely the specific forms of life may have changed. In catechetical instruction he will not proceed to the discussion of a new section without previous logical analysis; and he will not allow specific cognitions to remain isolated, but will determine their connection with previously gained cognitions and thus aim to weave all mental possessions into a unity — the unity of the Christian view of life. Thus religious instruction will gain in value to the pupil; his interest in it increases; he is securely fortified against many a doubt and against many of the inevitable attacks against his faith.

The most important forms in which logical thinking moves, or deposits its results, are \emph{conception}, \emph{judgment}, and \emph{conclusion}. The soul forms a \emph{conception} when it combines into a unity of thought the (essential, necessary) marks common to all individual objects belonging to one class or kind. Four steps are taken in this process: reproduction, reflection, abstraction, and combination. If, e.g., the \emph{conception} ‘tree’ is to be formed, the soul must first reproduce the images of all trees formerly seen; then it must carefully reflect what marks are common to all these trees and which of the marks are essential and which are unessential; it must separate the former from the latter; and finally it must combine all the essential marks into a well-rounded complete thought unit or \emph{conception}. The totality of all objects covered by one and the same \emph{conception} is called its extension, while the sum
\quellseite{238}
total of essential marks is called its content. By observing the extension we obtain the genus, and by observing the content we obtain the distinguishing marks; combining the two and expressing the result in terms, the soul renders a definition: the tree is a plant (genus) consisting of root, trunk, and branches (content of the conception, distinguishing marks). But by associating the two concepts ‘tree’ and ‘plant,’ and by fixing their relation to each other in the form of a sentence, the soul has progressed to the second stage of logical thinking: it has formed a judgment. It is the very purpose of judgment to indicate the relation of two concepts; asserting this relation in the form of a sentence or proposition we form a judgment. This, however, does not exhaust the reasoning powers of the soul. It does not only combine two concepts into a judgment; it can also unite two judgments and thus produce a third, a new judgment. In as much as we conclude something from known judgments, we name the resultant judgment conclusion. The soul draws conclusions in three ways: (1) by way of deduction, the specific being deduced from the general (all men are mortal—major premise: Caius is a man—minor premise; hence Caius is mortal—conclusion); (2) by the method of induction, the general truth being established on basis of the evidence of many specific cases (heat expands iron, it expands water, also air, etc.; hence it is concluded that it expands all objects); (3) by the method of analogy, the agreement of objects in two or more aspects suggesting and warranting their agreement in additional aspects (a hen protects her chicks, should not likewise a mother protect her children?).

Kant makes a further distinction between \emph{Verstand} (understanding) and \emph{Vernunft} (reason). \emph{Verstand} he defines as the soul’s faculty to deal with finite relations and dependent being, while \emph{Vernunft}, the highest faculty of the mind, according to him deals with infinite and independent being; it is the faculty of “ideas”. It is evident, however, that \emph{Vernunft} connotes no further advance in the soul’s activity, but
\quellseite{239}
merely a difference in regard to the object with which the same faculty is occupied. While it deals with the supersensible, the absolute, the infinite, the perfect, the divine, it is after all the same faculty that is in question, the same in its essence, forms, and laws. For this reason no further description of it is required.

In the activity of its thinking faculty the soul reaches the highest point of its intellectual life; but this is not to be understood as though rational thinking were possible only at the end of a long process of cultural development. Rational thinking rather is part of the pupil’s original and natural equipment, and often proves its existence by means of questions in the earliest years of life, just as the oak tree potentially is found in the little acorn which also betrays and asserts its individuality in the first years of its life. But as the young oak tree can develop its peculiar powers only where weather, soil, nurture, etc., are favorable, so also the thinking activity of the soul will be developed to the highest degree of efficiency only where sensations and concepts are normally formed and where instruction and training constantly exercise and improve this innate faculty.

The teacher of religion who realizes that he on his part has a share in the mental development of the pupils and that the material which he offers in religious instruction must be strongly fortified against doubt and temptation, can not possibly neglect teaching his pupils how to digest intellectually the religious material, how to form distinct conceptions, to render correct judgments and to draw logical conclusions. He must guard against passing on to the pupils ready-made conceptions; instead he will have them, on basis of vivid intuitions (from Biblical History as well as from daily life), gather the material themselves and discover all the marks of distinction; but he will not permit them to be overwhelmed by all the intuitions, but teach them to combine the discovered marks of distinction into conceptions. For as “conceptions without intuitions are barren,” so “intuitions
\quellseite{240}
without conceptions are blind.” The more central a conception is, and the more easily it is emptied of meaning or misinterpreted, the more thoroughly the teacher of religion will endeavor to obtain—his class cooperating—a precise and exact formulation and the more frequently he will return to it in the precise terms in which it has been fixed. He will teach his pupils to form independent judgments, never be satisfied with results obtained through guessing, never allow any gaps in the thinking process, but advance step by step in the development of a truth. He will lay particular stress upon the formulation and inculcation of clear moral judgments. For that reason he will when reflecting on actions of men and when criticising men, unceasingly point out the guiding motives and thus lead the children to the conviction that it is the motive which determines the worth of any action. Examples from contemporary life will prove particularly helpful in this respect. If conclusions are to be drawn, the path of induction will be primarily followed: from the example he will ascend to the rule, from the observation of specific instances to the law, from the special to the general; and he will not draw the conclusions himself, but have his pupils draw them. One of his main objects will be to unfold the truth by means of the question method and thereby to train the pupils for independent mental activity. If, finally, the teacher in partnership with the pupils has thoroughly treated and digested the material, then he will also require them to clothe the new truths in clear coherent language; the measure in which the pupils are able to do this, will be to the teacher a reliable criterion by which to gauge his teaching success.

\section{2. The Emotional Life}

The life of the soul is not limited to the exercise of the intellectual faculty (sensation, concept, thinking); it also has an emotional life. The processes of perception,
\quellseite{241}
concept, and thinking are accompanied by various conditions of pleasure and displeasure, such as sensuous pain and sensuous delight, joy and sadness, fear and hope, disgust and enthusiasm, delight in the beautiful and good, abhorrence of the ugly and hideous. These and similar excitations and states of the soul we call feelings or emotions. While they are never aroused apart from sensations or concepts or acts of thinking, they are essentially different from these. By the exercise of the intellect we learn that things are and what they are; through the emotions, however, we learn what \emph{value} things have for us. In proportion to their worth or worthlessness, things move our soul as soon as they have entered into our consciousness. Emotion, then, is the organ of evaluation; as such it is the source of every interest, a fact which explains the superlative importance which is attached to emotion.

The feelings of the soul may be classified according to their content and to their origin. According to \emph{content} they are either emotions of pleasure (e.g. the agreeable feeling of a pleasant taste or a fragrant smell, of a task successfully completed, the delight in the beautiful) or of displeasure (e.g. physical pain, hunger, anxiety, grief, discontent, ennui). According to \emph{origin} they are either sensory or mental. Mere sensations suffice to arouse sensory feelings; they are the concomitant phenomena of sensations. When, e.g., a piece of \emph{sugar} is taken into the mouth, the sensation then produced is accompanied by a feeling of pleasure. \emph{Sensory feelings} do not postulate any concept or thinking activity of the soul: they appear even in the earliest years of infancy. In the very first weeks of his life the child feels that it is agreeable to have enough food, to be warm, to hear a soothing lullaby, and likewise that it is disagreeable to be exposed to hunger, cold, or dazzling light. But the mental feelings postulate the faculty of the soul to form concepts
\quellseite{242}
and thoughts, and are awakened by these alone; they presuppose a certain degree of mental development. Since it is mental emotions that we are primarily concerned with, they require closer attention. We enumerate the following emotions: intellectual, esthetic, moral, religious, social, and self-feelings.

a. The search for truth is a trend innate in the soul. It feels genuine satisfaction, yes, joy and delight when after strenuous mental labor it has solved a problem or discovered the truth. Every enrichment of knowledge, every mental accomplishment, every new evidence of mental power and progress gives rise to such joy, is accompanied by such feeling of pleasure. But when the soul fails to make progress, when a problem must remain unsolved or is but partially solved, then it feels displeasure, discomfort, discontent. All such feelings are called \emph{intellectual feelings}, being emotions of the soul that accompany the success or failure of intellectual effort and register the value that we attach to success or failure. They are of the greatest significance for man’s mental development in that they prevent him from being content with false or immature results of mental efforts, they do not permit him permanently to give up the search for truth, but spur him on to further research until the correct solution shall have been found. Now, the teacher, also the teacher of religion, will most zealously cultivate these intellectual emotions; they are his best allies in teaching. He will never transmit to the pupils religious or moral truth as a finished product, but will rather enable them to feel the joy of seeking and finding and succeeding, and for that reason will engage jointly with them in the quest for truth, whether the material at hand be Biblical History or anything else. So conducted, his instruction will not be langweilig, it will not bore or cause ennui; the pupils will enjoy the privilege of cooperation, their interest will be awakened
\quellseite{243}
and sustained, the material becomes dear to them and develops into a source of joy. Aware of their powers, every effort crowned with success encourages them to undertake more difficult tasks; and joy over past achievements vouchsafes to them similar experiences as they make further progress in their quest for truth. Hence the teacher should never be satisfied with incomplete or partly correct answers, but lead the pupil to full comprehension of the truth; conversely, he will never reject a partially correct answer as though it were entirely wrong, but accept it as far as it is correct, thus gratifying the pupil and encouraging him to continue the search until the true answer be found. Many a pupil, especially in our country, does not learn to think and does not feel the desire for further mental work until he receives religious instruction; but it is also true that many a teacher of religion has shaped his instruction in such a manner that he imposes no mental tasks upon his pupils, never engages their cooperation, and as a result fails to gain their interest, to arouse their joy, and to achieve permanent positive results.

b. Another sense peculiar to the soul is the sense of the beautiful. A beautiful landscape, beautiful music, a beautiful poem well recited touches the soul most agreeably, arouses warm delight; but by the ugly and hideous it is displeased, repelled, disgusted. Such emotions aroused in the soul by the beautiful and the ugly are called \emph{esthetic emotions} (\emph{αἰσθάνεσθαι}). When we observe regular, symmetrical proportions and features, a certain harmony pervading every part of the form, we are attracted and pleased, especially so when the form is filled with precious content; but in all its fulness we behold the beautiful and feel its attraction when form and content correspond and blend perfectly and produce unity of effect, when the form, as it were, represents the embodiment of the ideal, the supersensible, the
\quellseite{244}
perfect, the eternal. Cultivation of the esthetic emotions enriches and ennobles the inner life, safeguards it against rude and vile influences and erects a barrier that may prove very valuable against wickedness and immorality. The teacher of religion will therefore not fail to cultivate these emotions. The wonders of creation in its vastest and in its minutest aspects, similarly the image of Christ and of many of his followers, or the image of man (1. Article), may aptly be utilized by the teacher for awakening or nourishing this sense of the beautiful. He will take time now and then to rejoice, with his pupils, in some product of sacred art, a beautiful picture, a beautiful melody, or a beautiful poem. And he will endeavor to keep away and combat everything vile and ugly in order that his pupils’ displeasure at it will increase in intensity to abhorrence and loathing.

c. Still higher in the scale of emotions are the \emph{moral feelings}. Abraham’s unselfishness and peacefulness, Jonathan’s friendship, Stephen’s courageous confession, arouse in our souls joy and pleasure. But Ham’s impiety, Absalom’s rebellion, Judas’ betrayal of the Lord, arouse in us displeasure and repugnance. Such emotions we call moral inasmuch as they indicate pleasure in what is morally good and displeasure in what is morally objectionable. The soul, in being aroused and pleased or delighted by what is morally good, registers the value that the morally good possesses for us. Being filled with displeasure at sin, the soul reveals the futility and worthlessness, yes, the hurtfulness and dangerous character of that which is morally objectionable. The moral feelings are, accordingly, valuable safeguards against wrong-doing and effectual guides and inducements to well-doing on our part. Here is found their eminent significance for the whole development of our inner and outer life. Among their outstanding results we mention the sense of duty, the sense of justice, and the sense of truth and honesty. At first
\quellseite{245}
generally unclear and elusive, they can be made keen, resolute, and strong by training so that, under their influence, clear moral images and conceptions are formed and healthy moral judgments become possible. The sum total of moral feelings, conceptions, and judgments we call “\emph{moral consciousness}.”

One of the primary tasks of the teacher of religion will be to deepen and to strengthen the moral feelings in the pupil’s soul, to form clear-cut moral images and conceptions, and to train him to form healthy moral judgments. The pupil as he is constituted by nature has, indeed, some understanding for that which the natural man’s morality considers good and evil, though even this knowledge is often defective; but he does not know, in the state of depravity, and needs to be taught, what is really good and evil in the sight of God. When a certain act is to be evaluated, the teacher will induce the pupil to consider the act as such, apart from its consequences. In this way he will train him to understand that a good act can never become evil on account of possible detrimental consequences and that an evil act can never be sanctified by eventual beneficial results, that every act is freighted with its moral value. The more clearly the pupil understands this principle, the more securely he will be fortified against the conventional utilitarian view. Furthermore, the teacher will train the pupil to penetrate from the overt act to the motive behind it and to recognize that here the worth or worthlessness of any act is determined. He will untiringly emphasize that nothing is truly good save that which spontaneously grows forth from fear and love of God or, to name the ultimate cause, from faith. He will employ every available instance from Biblical or Church History which may possibly help in producing keen moral convictions, in setting forth the value of noble moral character. He will guard against passing unbalanced or unsound moral judgments and endeavor to set an example of true
\quellseite{246}
Christian morality. But he has not even begun to appreciate his task if he has merely trained the intellect of the pupil to know right and wrong or if, in a legalistic manner, he exemplifies his precepts in his own conduct. He has influenced his pupil’s intellect, undoubtedly; it has received and repeated ready-made judgments and formulas. But the soul has remained cold, has not been stirred or thrilled. Yet, only that is important for one’s life, only that moves the will to action which has aroused the interest of the soul and stirred the emotions, which has awakened joy or pain in the soul and thus has received proper evaluation. “Values cannot be known as values without the \emph{inward experience} of them.” On the one hand, the teacher’s aim must be, therefore, to present moral goodness as it is evidenced in the lives of concrete personalities, with the utmost degree of sympathy and inner appreciation in order that the pupil’s soul may be attracted and stirred to genuine delight. Only then will moral excellency appear to him as valuable and worthy of imitation. What is morally evil, on the other hand, should be pictured with moral earnestness and unsimulated loathing; the pupil should see evil in all its ugliness and hatefulness as it starves the soul and poisons the wellsprings of life, so that he will catch something of the teacher’s scorn and loathing. Finally, the moral life of the teacher dare not bear the stamp of legalism and constraint; it should be free and gladsome, impressing the pupil how true morality elevates, emancipates, heartens, gladdens, and blesses (James 1:25). Now the pupil’s soul will be moved and he will value highly, and strive after, true morality.

d. From the moral feelings are to be distinguished the religious feelings, that is those emotions which result from the thought of God and our relation to him. Even natural man experiences to some extent such religious feelings: he feels himself overcome by some higher power, he bows beneath
\quellseite{247}
its relentless arm, he trusts in its assistance or invokes it. To a much higher degree such feelings move the Christian who worships God as holy Love and knows him as his Father in Christ Jesus. The Christian’s thought of God can not but be accompanied by emotion; the soul is dismayed at the energy of his holiness which casts out the sinner from his presence, it is reassured and blessed when it thinks of the blessed communion with God which is ever open to it. Here we observe the soil in which the moral feelings must take root if they are to attain to greater vigor; for if a morally good disposition or a good deed in itself elicits a feeling of pleasure and proves its attractiveness and value, how much more delightful will this pleasure become when the soul realizes that the disposition or deed is good also when measured by God’s criterion and that it will stand the test of judgment divine! Feeling such genuine joy over moral motives and acts, the soul places a higher value on them; appreciates their abiding and exclusive value, desirable above all else; lets the divine approval inspire courage and hope to attain it; and thus sets the will in motion to engage in its worthy quest. It is clear that no stronger incentives, no more efficient levers for the will exist than the religious feelings. God granting the strength (Phil. 2:13), moral acts become possible.

The union of moral and religious consciousness we call \emph{conscience}. Conscience judges of moral acts in their relation to God and determines their worth and worthlessness; thus it incites us to do good and deters us from evil, and lets us taste, in blissful happiness or in tormenting pain, the worth or worthlessness of the moral act.

Religious instruction exists for the purpose of training the religious feelings. The teacher’s aim is to place the souls of his pupils face to face with God so that they will walk before him, measure everything by his standard, and find
\quellseite{248}
joy and happiness only in his fellowship. In order to achieve this aim it will be necessary that the teacher shows, at the hand of \emph{living examples}, what true religion is, how it makes humble and modest in prosperity, strong and composed in adversity, how it supplies life with its most precious contents and affords the only help in the agony of death, how it frees the soul from sin and all fears and worries and equips it with strength to battle against every form of impurity and ungodliness, how it is the sole fountain of true, of permanent happiness. In proportion as the pupil beholds, \emph{actually beholds}, the value of religion in the form of living, historic persons, and, in a measure at least, also in the life of his \emph{teacher}, he will stretch out his arms for it and endeavor to follow that path on which it promises to \emph{enter} into his life. When our heavenly Father sent into the world his only Son to become flesh, was not one of his purposes this that we should behold in him, incarnate, both His holiness and His love so that our soul might be stirred and we might look up to Him, the only One whose pardon, whose favor is the sole object of consequence in heaven and on earth? One who, during religious instruction, has never seen and tasted how friendly and kind God is, may know ever so much about religion—a vital and desirable reality it has not yet become to him.

e. As the fifth group of emotions we mention \emph{social feelings}, or the concomitant phenomena of our fellowship with others. Most important among them are the sympathetic feelings, which are caused by perceiving or thinking of the weal and woe of our fellowman and which let us share his gladness and sadness. When the Samaritan knelt beside the Jew and assisted him or when Jesus wept in view of Jerusalem’s impending doom, social feelings had been aroused; and they were stifled by Shimei when he cursed David, and by the priest and the Levite when they passed
\quellseite{249}
the man fallen among the thieves. Man is enbruted when these feelings are not cultivated: coldly he ignores his neighbor’s grief or joy, and he himself is the sole object of his thoughts. Where they are cultivated, however, their response increases ever more in swiftness and vigor, and they arouse to energetic help in case of suffering. How can the teacher of religion expect his pupil to fulfil the duties of neighborly love, to associate with others in common tasks, to participate in common worship, to cooperate in the sphere of foreign and inner missions, if he does not, in instruction, zealously cultivate and strengthen the innate sympathetic feelings of his pupil? While the family is the \emph{primary nursery} of the social feelings, the teacher of religion will not fail to \emph{quicken and strengthen} them. He will observe his pupils’ conduct toward one another; he will relentlessly combat cruelty, envy, or malice, and attempt, whenever such traits manifest themselves, to quicken the social feelings. He will scrutinize every part of his educational material for suitable object-lessons (e.g., how Abraham delivered Lot or how he interceded for the cities, the story of the Young Man at Nain, the parables of Luke 15). Above all he will find abundant pertinent material in the life of Jesus “who had compassion on the multitudes because they were as sheep not having a shepherd.” Of profound importance is the teacher’s own example: the teacher who always has time to enter into the joys and griefs of his pupils with understanding and sympathy, will almost certainly observe how their souls vibrate when when they come in contact with the grief and glee of others; they have experienced the soothing effect of tender sympathy and know its transcendent value. — Among the social feelings must be numbered also love, respect, confidence, and gratitude. These feelings, likewise, express the value that another person possesses in our estimation. The teacher, needless to say, has duties to perform also in the strengthening of these emotions.
\quellseite{250}

f. The soul, finally, may be moved when it considers one’s own person, its worth or worthlessness. These feelings are \emph{self-esteem, sense of honor, and shame}. They are by no means unimportant for the teacher of religion. He will recognize the propriety of self-esteem and endeavor to arouse it in those who are inclined to be timid; where there is no self-esteem or self-respect, a courageous moral act is unthinkable, and the battle of life, particularly the moral battle, is abandoned before it is fairly begun. More frequently, however, the teacher will have to check and curb the sense of self-esteem lest it develop into conceit, pride, and vanity. The same is true of the sense of honor. A person who has no sense of honor and is indifferent to praise and censure, exposes himself to the danger of moral indifference; but when this sense is excessively developed, it degenerates into offensive sensitiveness where one continuously complains of being slighted or neglected and makes life unbearably miserable for others. Special attention should be given to the sense of shame not only as the guardian against many a sin, but also as the morning dawn of a better life if sin has been committed. True repentance does not exist without shame for one’s own moral failure, without anguish over the wrong committed or over the ugly and mean phases of one’s conduct. But mere indoctrination will not suffice; sin must be viewed at the hand of living examples, its baseness must be clearly apprehended if the pupil is ever to recognize, feel, and experience his own sin as failure and guilt for which he would hide from both God and men.

So rich and varied is the soul’s emotional life that we may speak of intellectual, æsthetic, religious, and other feelings. And yet, what is really experienced by the soul, is invariably either pleasure or displeasure, everything else describing merely the condition or sphere upon which the soul reacts in such a way as to bring about either joy or
\quellseite{251}
grief. Again, however numerous the conditions or spheres may be in regard to which the soul is stirred, the activity of the intellect is always postulated. Not until objects or conditions have entered into consciousness, can the soul be stirred by them: the emotions are concomitant phenomena of intellectual processes. \emph{Important concomitants, indeed}; for they are the criteria of the value of the objects of the intellectual world. It is true, they are not absolute and objective criteria inasmuch as the conception, and the value, of the moral good are by no means commensurate with the degree of pleasure aroused by the good in the soul; and guilt remains guilt whether one feels it as such or not. Moreover, the degree of the development and training of these emotions, to some extent also the congenital temperamental tendencies, cause one and the same object to arouse one person more intensely or speedily than another. The response of the sanguine temperament is swift and feeble, that of the choleric temperament is swift and powerful, that of the phlegmatic is slow and weak, and that of the melancholy temperament, while slow, is strong and deep. Although not objective and absolute criteria, the emotions are nevertheless subjective criteria of momentous significance inasmuch as they permit us, in a measure, to appreciate and to experience the value of things and objects. If the soul did not possess the power of feeling, it would be devoid of an organ with which to measure the value of the objects of matter and intellect; it could not possibly experience their value however valuable the objects might be and however frequently the soul might have—purely intellectually—endorsed the value judgments passed by others.

The emotions considered so far are excitations aroused in the soul in the usual order of events. Emotions caused by unexpected stimuli and manifesting themselves with considerable intensity, though generally only for a short period
\quellseite{252}
of time, are called \emph{affections}. Owing to peculiar circumstances, already existing feelings may be so intensified as to develop into \emph{affections}. Intellectual emotions, e.g., may develop into embarrassment, bewilderment, astonishment, enthusiasm; esthetic emotions into admiration, rapture, transport; moral and religious emotions into indignation, shame, remorse, ecstasy; the social feelings into fervent love, passionate hatred, rage; the self-feelings into arrogance, despondency, fear, despair. More permanent emotional states whose origins are not clearly recognized we call \emph{moods}. Affections as well as moods can and should always \emph{remain under control}. The teacher of religion will most successfully lead his pupils to exercise such control, if he makes his pupils feel, at the hand of living examples, how miserable are the effects of most of them, how wretched men make themselves by yielding to them, and how they debase themselves to the rank of the slave while they ought to be masters of their affections. In such demonstrations the teacher must often corroborate his lessons by drawing upon the moral and religious feelings. The sum total of emotional phenomena we call \emph{temper}, \emph{Gemüt}.

Most closely connected with the emotional life is what in psychology and pedagogy is called \emph{interest}. It is the opposite of the mental state called \emph{indifference}. If my soul is indifferent to an object, it reacts or vibrates to it in no wise, it feels no pleasure or displeasure, its presence is no gain to me and its absence no loss. But if my soul is interested in something, it can not but be of value to me, it stirs my soul to gladness so that, delighted with it, my soul prefers it to other objects and gladly occupies itself with it. The object has entered into a vital relation to me, a sort of reciprocity has been established between it and my soul, in brief, it has become valuable to me and I have become conscious of such value. Now it is evident that interest is closely
\quellseite{253}
connected with the emotions, for, as was seen above, the emotions constitute the soul’s organ of evaluation. Hence we may speak of sensuous interests and mental interests, and enumerate among the latter such as intellectual, esthetic, religious, or social, interests. Yet, the difference between interest and feeling dare not be overlooked. On the one hand, while my feelings may indicate the worthlessness of an object which has caused pain or displeasure to me, the term interest is usually employed in the positive sense, viz., whenever an object elicits our pleasure and proves its value to us. On the other hand, the organ of evaluation must not be confounded with the resultant state in which the soul prefers that object to others and gladly occupies itself with it. But just that state is interest. In the measure in which my soul has been stirred, time and again, to joy by an object or by recalling the joyful experiences previously occasioned by it, my interest in the object in question is deepened and strengthened. My evaluation of the object will become more pronounced as it continues to affect my soul in ever larger measure, and it will be assured of my permanent interest if its effects upon me remain exclusively, or largely, beneficial. — If the emotions were correctly represented as the key to the will, the soul striving after nothing which the emotions have not first declared to be desirable, the same is true, in an even higher degree, of interest. Interest is the real motive power of all endeavor and volition. Certainly the teacher of religion will not neglect its importance since he—more than other teachers—purposes to influence the will of his pupil to action.

\section{3. The Volitional Life}

The third element of the life of the soul is its volitional life. The soul does not only become conscious of the objects of the world without and within itself; it does not only register
\quellseite{254}
the value which it attaches to these objects; but the soul draws, from the knowledge it has obtained and from the evaluation placed upon the objects, a practical conclusion: it stretches forth toward that which is precious or attractive in order to grasp it, and it spurns that which is unattractive and displeasing. This \emph{striving} for objects conceived and appraised by it, constitutes the volitional life of the soul. Volition is by no means merely a derivative product of, or some faculty identical with, our intellect or our feelings, but a faculty of the soul as original and real as the other two. While the intellect and the emotions deal with the objects of the present moment, the soul, in its volitional life, turns to the future with the purpose of effecting a change in the circumstances obtaining in the present, either as regards its own state or that of the outer world as conceived by it. In view of the fact that the terms “volition” and “will” may be taken in a variety of meanings and in their narrower sense denote the highest aspect of the whole volitional life of the soul, it would be better to use the term “striving” for this phase of soul life. The striving of the soul proceeds in three stages: impulse, desire, and volition.

a. Striving is felt in the initial stages of human life, even before concepts are formed and experiences are gained. Resting in the womb of its mother, the child strives and presses outward although the outer world has not been perceived or experienced. The new born babe cries for food for the purpose of self-preservation although it has neither seen food nor experienced that food will allay hunger and bring about the pleasant feeling of satiety. When it has lain quiet for a while, it craves motion and kicks until it has freed itself although it never experienced freedom of motion and its agreeable effects. Without ever witnessing sociability or experiencing its value, it longs to be near its mother and later with other children. There dwells in the soul a hidden force, an unconscious striving which stretches forth after various things. This is volitional
\quellseite{255}
life in its lowest form, the \emph{impulse}. Since impulse manifests itself in the mental no less than in the sensuous sphere, it has become customary to speak of sensuous and mental impulses although this usage is not quite free from objection. To the former class belong the impulses of nutrition, self-preservation, and propagation; to the latter those of sociability, imitation, purposive activity, and acquisition of knowledge. Unconsciously, in every case, a certain power asserts itself that makes for the gratification of important needs. If the \emph{desire} is not satisfied, body and soul may be stirred to positive pain; and if the desire is satisfied, body and soul are set at rest, at least until the impulse in question once more begins to stir, and the longing of the soul for this or that object again manifests itself: for impulse is something \emph{permanent}, existing for the whole space of life or, at any rate, for the space of a certain stage of life, and manifesting itself again and again in the same strivings. In the sphere of animal life instinct (\emph{instinctus}=impulse) corresponds to impulse. No essential difference would be noticeable between man and animal in this respect if human volition had no further stages. This implies that the impulse, originally, is neither good nor bad; it is, as a matter of fact, morally neutral, being creatively an essential part of human nature. But it is equally certain that in the state of sin the impulse tends to exceed its bounds; it must needs be controlled, a fact which suggests that the teacher of religion has a duty also in this direction.

b. The impulse as such is blind; it lacks a clearly recognized aim. When, however, upon repeated gratification of the impulse, the soul has learned the means of such gratification, then the impulse assumes the more definite form of desire. Striving no longer is vague or purposeless, now it has an aim, it would attain some object concerning which experience has taught that it affords gratification of the impulse and the feeling of pleasure. In distinction from the impulse, it is seen, desire
\quellseite{256}
is brought about upon the ground of intellectual and emotional activities of the soul. Now the objects conceived by the intellect are variously evaluated by the emotions, some of them as valuable because agreeable, and others as detrimental because disagreeable; hence desire represents now a positive, now a negative striving of the soul, or in other words, now it is a striving for the possession of the object in question, now it is aversion to it. Desire is usually of transient nature: how rapidly its objects shift in the case of the child! But it may assume a more vigorous form and thus attain to permanency. The scale it may possibly traverse is given in the words: desire proper, \emph{inclination}, \emph{propensity}, \emph{passion}. \emph{Desire proper} is an intensive desire, the soul striving to seize an object with energy and vehemence. This vehemence may disappear as soon as the desire has been gratified as is seen in the case of the thirsty person as soon as he has slaked his thirst. But it may continue returning, and so become a habit, as in the case of the drunkard. In consequence of repeated gratification desire increases in strength just as the repetition of one and the same concept renders it clearer and more vigorous (cf. p. 207). Desire having become habitual is fitly called inclination (to be distinguished from native tendencies) because the trend of the soul's striving at this stage always points in a definite direction. When inclination has attained to a certain degree of intensity, it has become propensity: as the tree finally leans altogether toward the side (\emph{propendere}) to which the wind has continually blown its branches, so eventually the whole trend of life turns to the object toward which one has been drifting with increasing frequency and force. While propensity is already virtual enslavement, passion is a desire of such strength that it seeks gratification in a certain direction, no matter what the cost; reason may condemn it; conscience may protest against it; its consequences may be known to be ever so wretched. It is true, the terms “propensity” and “passion” are used also
\quellseite{257}
to characterize the desire for the morally good; but this usage is by no means prevalent. This is not true of the words “desire” and “inclination”: these are neutral and receive their ethical color from the object to which they are directed.

If already the impulses of the pupil occasion special tasks for the teacher of religion, this is true in an even higher degree of desire in its several degrees of intensity. The teacher will see on his part that justified sensuous desires are gratified, thereby promoting the health of the body and furnishing the basis for genuine child-like cheerfulness. Undiscriminating and insatiable desires, however, he will try to inhibit by occasional admonitions, also by his own example, and by pointing out the value of noble simplicity and self-restraint. He will pay close attention to his pupils’ conduct, to the friendships they form, to the books they read, to the pictures they see, and he will combat every form of unchastity and lewdness in words or gestures lest the dormant sexual impulse be prematurely aroused or sexual desire when already awake be fed; and he will watch over his pupils lest any tendencies toward indolence, busybodiness, seclusion, tyranny, greediness (e. g. in games), gossip, prejudice, etc., should be established. He will direct the pupils’ desire especially toward goodness and godliness; and by arousing interest therein, produce and strengthen the desire and striving to attain those. Conscious of the great power of habit, the teacher will make an effort to surround all of his instructions with holy exercises, to begin with prayer and to close with prayer; to investigate the attitude of the home toward prayer; to urge regular attendance upon public worship; to pay attention to the exercise of charity in such forms as are adapted to his pupils, in order gradually to effect a certain degree of facility in well-doing — this characteristic of Christian “virtue.”

c. Higher in the scale of striving than desire is \emph{volition} (in the narrower sense). It postulates not only the knowledge
\quellseite{258}
of the object desired, but also a knowledge of the ways and means that lead to its acquisition, and, likewise, confidence in the adequacy of one’s powers to acquire it. The soul \emph{wills} when on basis of such knowledge and confidence it resolves to appropriate the object in question. The prodigal son, in his misery, may often have thought of his father’s house with its abundance, and more than once he may have felt the desire to return, but here the process ended. Not until one day when once more the image of his home most vividly gripped him, when he reflected upon the way that leads to it, and deliberated on the words which he should address to his father, when he felt confident that his father would not cast him away, when he overcame all doubts as to his own ability to carry out the proposed plan, and then decided to venture upon the chosen course—not until then can it be said that he willed.

Volition, it is seen, postulates the activity of intellect and emotions even more sweepingly than does desire; volition becomes possible only through vigorous activity of both intellect and emotions. Corresponding to Kant’s above mentioned distinction between \emph{Verstand} and \emph{Vernunft} (cf. p. 216), the distinction is made between intellectual and rational volition of the soul. The former has for its object the useful and expedient, the latter the moral good and divinely acceptable. The former is determined, in its choices and resolutions, by advantages that may possibly be gained or losses to be avoided, the latter by the consideration of the mutual duties and rights of men and by the reference to the divine will, i.e. by moral and religious values.

As soon as volition exists, it urges the execution of the choice in action. The activity of the will is inward in that it focuses our attention on an object or reproduces and associates concepts or, chiefly, controls the current of our thoughts through reflection, i.e., voluntary meditation. Its
\quellseite{259}
activity is outward in that it leads to actions; for action does not denote any kind of deed or act, least of all an act which is the incidental response to some blind impulse, but such an act as takes place in consequence of deliberation and evaluation (thinking and feeling) and of a definite choice traceable to both of these (volition), or, in other words, an act that is performed in perfect freedom. For that there is a \emph{freedom of will} is plainly seen when we observe that the soul may decide on a certain course of action contrary to the definite knowledge that it is wrong and in spite of the fact that the emotions have appraised this course as worthless or even detrimental, or when we observe that the voice of reason (intellect) and the excitation of the feelings, be they ever so potent, by no means constitute coercion which the soul could not escape but must obey.

When a person repeatedly faces the same problem in moral and religious matters and repeatedly makes the same choice or decision, then such decision develops into what is called a maxim, i.e., a moral principle, a moral rule of life. Steadily conforming in all decisions and actions to this maxim, a corresponding stamp will be impressed upon his will, his conduct, his whole personality. Inasmuch as this stamp or mark constitutes a distinctive trait, an “impress”, of the person, we appropriately call it a characteristic, for character (from \emph{χαράσσειν}=to make sharp or pointed, to engrave, to impress) denotes the impress or stamp. When the soul is guided by certain maxims not only in one phase but in all moral and religious decisions, man has become a character. While this, in a certain sense, applies also to one who has become so confirmed in wickedness by repeated wrong-doing that he consistently acts in conformity to wrong principles, usage applies the term character preferably to such as act in conformity to principles morally good, whether we have in mind the principles of natural, or of Christian, ethics.
\quellseite{260}

If it is agreed that moral conduct and the development of character are possible only upon the basis of volition, the transcendent importance of the training of the will in education, especially in religious education, is readily seen. To find the way to the will of the pupil and to move it to action, must be the main object of the teacher of religion if he has learned that religion is not a mere aggregate of doctrines but rather an actual condition, namely: fellowship with God through Christ and the actual recognition and execution in life of the divine will. He will most frequently impress upon himself the fact that the will can be touched only through the intellect and the feelings; that, accordingly, his aim must be directed, indeed, at the formation of clear conceptions and the gain of accurate knowledge, but that this process should not take place in a dead or formal manner, that it must rather be accompanied unintermittently by warmest interest and by the highest regard for the truths in question. Thus the feelings of the pupil, agreeably touched, will vibrate to the teacher’s treatment; his interest in the truths under consideration will be aroused; he will begin to share his teacher’s evaluation of these truths. When, moreover, through representatives of a noble morality, and the teacher’s own personality, the fact is brought home to the pupil that the moral good is attainable, and when also the way to such attainment is shown, then the pupil does not merely feel a transient desire for the moral good, but he actually wills it, decides to do it. If, however, the teacher would safeguard himself against disheartening disappointment, he will, from the very beginning, reckon with his own inability to force the right decision. The will despite its close relation to the intellect and to the feelings, is notwithstanding an independent factor. In the state of sin it is by no means always possible even to arouse the feelings in behalf of the good, in part because the direction of the will reacts upon the feelings and even
\quellseite{261}
upon the intellect, in part because through original sin these organs have been seriously enfeebled for the performance of that which is good. In the state of sin, we must remember, the “\emph{propensitas ad malum}” has replaced the receptivity of intellect, feelings, and will for that which is truly good and pleasing to God. And although by regeneration and justification a new life has been established, it is to be kept in mind that the old life is not extirpated but rather projects itself into the new and that, therefore, the response of the intellect, the feelings, and particularly of the will, is by no means so spontaneous or complete as may be desired and expected. Nevertheless, the teacher will not cease trying to influence the pupil’s will or the atmosphere of the pupil’s home; he will be encouraged in this by the well established facts of the greater receptivity and responsiveness of youth. He will not only induce his pupils to form decisions, but also seek to bring about corresponding actions. The spheres in which he will guide his pupils to right actions are the associations and friendships formed by them, their home life, their relation to pastor and teacher, and their participation in the work of the congregation. In the course of time the power of decision itself will be affected by such training. The teacher will remember the vast importance of repetition, of exercise in well-doing, of familiar acquaintance with it, and will thus, by purposeful habituation, lay the foundations of a Christian character.

Even this brief delineation of the life of the soul affords an insight into its wealth and manysidedness.

We have distinguished three fundamental powers of the soul, intellect, emotions, will; yet, the soul is a unity. It is the same soul that knows, feels, and wills. The soul is the subject of the threefold activity of knowing, feeling, and willing. It is, as it were, the center of the three spheres of vital activity which we call intellectual, emotional, and volitional
\quellseite{262}
life. It is conscious of the fact that these spheres are its own, clearly distinct from anything foreign. For this reason we speak of \emph{self-consciousness, or the consciousness of the Ego}. The soul can distinguish its activities from itself, can place them before its vision and make them objects of research; but what is thus objectivized, is always a part of itself. Furthermore, just as in forming the concept of a blooming tree the soul fuses into a \emph{unity} a variety of concepts, even so it fuses into a unity its own three vital activities. We, therefore, speak of the unity of the ego-consciousness. Finally, all organic changes notwithstanding, the soul is conscious of its former sensations, concepts, states, strivings, no less than of those of the present or the future. Hence we speak of the \emph{continuity and identity} of the ego-consciousness. It is especially the last named factor which renders absolutely untenable the view that the soul is no more than an activity or an attribute of matter (of the body or the brain which are subject to organic changes). The soul is a real, independent, self-subsistent entity, neither occasioned by matter nor disappearing with the decay of matter—however little we may be able to assert in detail concerning its ultimate nature.

\quellseite{263}
\chapter[§ 22. The Gradual Unfolding]{§ 22. The Gradual Unfolding of the Pupil’s Life}

Stanley Hall, \emph{Ausgew. Beitr. z. Kinderpsychologie u. Paedagogik}, gesammelt von Ufer, 1902 (These essays were published in The Forum, Scribner’s, Pedagogical Seminary, American Journal of Psychology, etc.). W. Preyer, \emph{Die Seele d. Kindes}, 1882, ¹908; English by H. W. Brown, 1889. K. Groos, \emph{Das Seelenleben des Kindes}, ²1908. Cl. and M. Stern, \emph{Monographien ueber d. seelische Entwicklung d. Kindes}, 1907 ff. G. Baeumer and L. Droescher, \emph{Von d. Kindesseele}, 1908. R. Gaupp, \emph{Die Psychologie d. Kindes}, ³1912. W. Ament, \emph{Die Seele des Kindes}, ⁴1914. Sigismund, \emph{Kind u. Welt}, ²1897. B. Otto, \emph{Beitr. z. Psychologie d. Unterrichts}, 1903. K. R. Loewe, \emph{Wie erziehe u. belehre ich m. Kind bis z. 6. Lebensjahr}, ²1904. R. Hartmann, \emph{Die Analyse d. kindl. Gedankenkreises}, ⁴1906. A. Schreiber, \emph{Das Buch v. Kind}, 1907. R. Goltz, \emph{Das Buch der Kindheit}, ⁵1909. G. Compayré, \emph{L’Evolution Intellectuelle et Morale de l’Enfant}; part i German by Ufer: \emph{Die Entwicklung d. Kinderseele}, 1900; part ii English by M. Wilson: \emph{Development of the Child in Later Infancy}, 1914. F. Tracy and I. Stimpfel, \emph{Psychology of Childhood}, 1884, ⁷1909. J. M. Baldwin, \emph{Mental Development in the Child and the Race}, 1895, ⁷1911. Fr. Froebel, \emph{Pedagogics.of the Kindergarten}, 1895. A. Kirkpatrick, \emph{Fundamentals of Child-study}, 1903, ¹⁴1914. W. Geo. Koons, \emph{The Child’s Religious Life}, 1903. Ch. W. Rishell, \emph{The Child as God’s Child}, 1904. G. H. Pease, \emph{An Outline of a Bible School Curriculum}, 1904, ²1909. Ant. Abern. Lamoreau, \emph{The Unfolding Life}, 1907. G. E. Dawson, \emph{The Child and His Religion}, 1909. E. P. St. John, \emph{Child Nature and Child Nurture}, 1911. H. Bauer, \emph{Die Psychologie der Jugendlichen}, 1911. J. Eger und L. Heitmann, \emph{Die Entwicklungsjahre}, esp. Heft vii: F. Mahling, \emph{Die Psychologie d. Jugendlichen}, 1913. G. Fuellkrug, \emph{Zur Seelenkunde d. weibl. Jugend}, 1913. F. W. Foerster, \emph{Jugendlehre}, 1904, ⁵1911. F. W. Forbush, \emph{The Boy Problem}, 1911, ⁷1913. John L. Alexander, \emph{The Boy and the Sunday School}, 1910. W. Buck, \emph{Boys’ Self-governing Clubs}. E. H. Hughes, \emph{A Boy’s Religion}. Fr. W. Johnson, \emph{Problems of Boyhood}. W. B. Forbush, \emph{Church Work With Boys}, 1910. Stanley Hall, \emph{Adolescence, its Psychology and its Relation to Physiology, Anthropology, Sociology, Sex, Crime, Religion, and Education}, 1904, ²1911. St. Hall, \emph{Youth. Its Education, Regimen, and Hygiene}, 1904, ⁴1911. T. King, \emph{The High School Age}. Marg. Slattery, \emph{The Girl in Her Teens}. M. E. Moxley, \emph{Girlhood and Character}, 1916.

\quellseite{264}
The catechist needs to know the soul life of his pupils in general, but no less essential is a knowledge of its gradual unfolding. Otherwise he might possibly accord his pupils a treatment altogether unadapted to the stage of their development. The child comes into the world a perfect human being, indeed, and his organism stands in no need of completion whatsoever; but he is as yet totally undeveloped. All the possibilities and powers of the fullgrown man or woman are found in the child, but only germinally, potentially; a development lasting from 18 to 21 years, at times even longer, is needed, before the human organism which has come into the world at birth, has fully matured. The observation of this gradual development is one of the most fascinating studies. Three stages are distinguished: (1) Early and later infancy, which extends from the first to the sixth year; (2) The childhood stage, which extends from the sixth to the twelfth or thirteenth year, and which is divided into two sections (six to nine, and nine to twelve); (3) Early and later adolescence, extending from the twelfth or thirteenth year to the period of full maturity. These stages are not definitely marked, and vary by a full year or more in different cases; but their existence can not be doubted and in general they correspond to the age given above.

\section{1. Early and later infancy}

The most interesting and the most exhaustively investigated of these stages is early and later infancy. At this place, however, we must confine ourselves to the essentials because the infant becomes a subject of instruction only to a very limited degree.

a.

When the child comes into the world, he finds himself suddenly surrounded by a gigantic chaos, full of light and sound, “a big, blooming, buzzing confusion,” to use the expression of William James. A number of stimuli enter upon the soul from the outside world, acting upon the perceptive faculty in such a manner as to call forth all sorts
\quellseite{265}
of sensations. These sensations are at first the only manifestations of the intellectual life. While the sense impressions are at first indistinct and undefined, they enlarge, already during the first year, into clear perceptions (the look of the child is fixed upon the shining watch from which a stimulus has proceeded; his eyes follow a light carried before him). Simultaneously we observe the first beginnings of concepts and the reproduction of such concepts (he has repeatedly heard his father say ‘tictac’ when handling his watch, associates the two concepts, and eventually points out the watch whenever he hears the sound ‘tictac’), of attention and of thinking (cf. pp. 213; 217). Even linguistic development generally sets in as early as the first year: the child begins to exercise his organs of speech by forming all kinds of spontaneous sounds with which, however, he connects no meaning as yet (‘tata’). The feelings of the child are at first exclusively of the sensuous nature: displeasure is aroused by hunger and cold, and pleasure by warmth and the gratification of the appetite (cf. p. 219). Still, toward the end of the first year of life faint beginnings of higher feelings appear: delight in color, fear of strangers and of unusual sounds, love for the mother. \emph{Volition} is expressed at first only in the form of impulses, especially those of nutrition and motion; but as soon as perceptions, concepts, and the reproduction of these begin to show themselves, there is also conscious desire (the child catches at the milk bottle, and cries when it is withheld).

b. A tremendous advance is made by the child when, in the second year, he begins to walk and to talk. In consequence of his ability and skill in walking, he not only conquers his helplessness more and more, but also obtains the necessary space, constantly enlarging, for new observations, experiences, and activities. The impulse of motion, active already in early infancy, now comes into its own; the power
\quellseite{266}
inherent in the child, later to be utilized in mental effort, now rushes outward in the sphere of bodily life with such energy that restlessness becomes fairly characteristic of it. Closely connected with the progress of bodily development is a significant unfolding of the soul life, chiefly of the intellect. The organs of sense, especially eyes and ears, now hungrily stretch out toward all attractions of the new and exceedingly interesting world, transmitting them, as wide-open canals, to the nervous system and thereby to the brain, which impressionable like delicate clay, receives and communicates them to the soul. The construction of the inner world sets in which will be poor or rich in exact proportion to the impressions working upon the soul from without. In his speech (language is dependent on imitation and inherited equipment) he gives expression to this world with ever increasing facility. Presently also \emph{memory} develops: the brain cells, exceedingly receptive during childhood, retain the impressions received; and as soon as the soul acts upon them, they are revived and the images formed in the past appear to the child anew. If the first impressions are made at a time when the brain is not weary and the child is attentive; if the object to be brought to the child's attention is made clear and impressive; if, in addition to the ear, also the eye (picture-book with verses), or some other organ, is employed; if the new is coupled with something already known and repeated one time after another; then \emph{memory} will prove quite faithful already from the third year. \emph{The phantasy} of the child is based upon \emph{memory}. This faculty awakens in the third or, more generally, in the fourth year, and at once begins to unfold a vigorous life. That is the reason why the child is not only fond of hearing tales at this age, but also constructs a wondrous world of his own out of the concepts stored in his soul; why at play he converts old objects into ever new and changing forms; and
\quellseite{267}
why he tries his hand at drawing. In consequence of imaginative concepts, akin to phantasy, the child often entertains fears where there is no occasion for them. As yet \emph{reasoning} is submerged under the sense perceptions and the play of phantasy which are far remote from exact logic. Notwithstanding it asserts itself toward the end of this period in numerous childlike investigations and experiments, and especially in innumerable questions in regard to the “what,” “wherefore,” “whereto.” — As far as the \emph{emotional life} in later infancy is concerned, the sensuous feelings which were thoroughly dominant during the first year of life, remain in the foreground also during the present period. It may be observed, however, how the intellectual emotions gradually unfold, and even the first promptings of the sense of honor, of moral, religious, and social feelings become manifest. But this is the case only when education does its part; for, left to himself, the little child is totally self-centered — instinctively so before he learns to distinguish himself from the beings surrounding him, and consciously so when he has attained to self-consciousness (third or fourth year). — Also in the sphere of \emph{volition} the sensuous preponderates. But, with the nobler emotional promptings already mentioned, also higher strivings begin to appear. The imitative impulse which is quite strong in the child in this period may prove of particular importance for the reason that, by its aid, the first habits may be formed which are of vast consequence for one’s whole life.

Whether and, if so, how the soul life of the child may be influenced at this stage, are questions to which widely divergent answers have been given. As far as we are concerned, we are certain — and modern psychology corroborates rather than invalidates this view — that the mother may exert pre-natal influence upon her child, that the mental and spiritual atmosphere in which the mother moves during the
\quellseite{268}
period of gestation does not leave the child without a trace. Unquestionably such influence is effective during the first year. While this applies to the unbaptized child, it applies even more strongly, as experience shows, to the baptized child, in whom the Spirit of God is dwelling and to all of whose soul-functions He has access. Thus influences from without if of the proper nature will meet the longing of the soul of the baptized child. It is self-evident that in the first years of life conduct can not be controlled by command and precept; in fact these should but rarely be applied even from the fourth to the sixth year. At this early stage, the chief factor of training is the moral and religious environment, the moral and religious atmosphere in which the child grows up. Silent and gentle as dew from heaven, and equally effective, is the way in which home life influences the slowly unfolding soul of the child. Insatiably the sense-organs stretch forth toward all the stimuli of the environment, extending their hospitality no less to those which are morally evil and worthless than to the good and wholesome ones; and whatever has once been received, is retained, and is bound, sooner or later, to exert its power. Memory being of such phenomenal strength, it will be the paramount duty of the parents most carefully to watch over their gestures, words, and acts, over the spirit governing their house. Cheerful fellowship coupled with faithful performance of duty and consecrated by regular family devotions can not but leave its traces in the soul of the little one. As early as the third year the child may be impressed, by the aid of pictures, with the permanent presence of God than which there is no stronger foundation for the moral and religious life. One of the most valuable allies of training in this respect is the impulse of imitation: observing his parents stand in the presence of God and speak with Him, the living and present One, the child will of his own account likewise fold his hands, speak to Him, address
\quellseite{269}
to Him all his requests, both large and small. This is, indeed, not conscious saving faith, but it is the relation to God which corresponds to the present stage of the child’s development and it will have lasting effects upon the entire subsequent life of the child. The child’s experience of God will be widened and invigorated if the parents will nourish the phantasy of the child which begins its activity about the fourth year; if—as Luther did so exquisitely in his letter to little Haenschen—they will present to him in most vivid colors an attractive picture of the beauties of the heavenly garden; if, with constant adaptation to the child’s limited horizon, they tell him gripping stories from the Bible or from every day life which show how the Father in heaven cares for all His children; or if they tell him of the bountiful goodness of God when spring clothes the fields with radiant verdure, and of the omnipotence of God when violent crashes of thunder frighten the child. The little ones ask countless questions; instead of suppressing this desire, parents will most warmly welcome it, not only because it indicates the first movements of reasoning activity, but especially because it affords them an opportunity to direct their children to God and to His works. When, e.g., the child asks about the new life that begins to throb with the arrival of spring, the parent should seize this opportunity to tell him about the creation of the world, of course, in language intelligible to the child. Or when he asks a question about the picture of the Crucified that adorns the wall, he will utilize it as the starting point for an account of Jesus the Saviour. It will be advisable, owing to the power of the child’s memory, to connect with such an account a rhyme or a hymn stanza and repeatedly to sing such stanzas with the children for the purpose of imbedding them more firmly in their memory.

The parents must recognize especially the importance of the emotional life of the child. Since only through the stimulation
\quellseite{270}
of the emotions, the way to the will is found and the soul is moved to action, the parents will endeavor to bring the proper influence to bear upon their children. They will achieve this result not so much through precept and admonition, as through suggestion, sympathetic radiation; for the intellectual life of most children is inadequately developed to react appreciably to mere commands, and, after all, whatever may be accomplished in this manner will rarely become a permanent trait of character, but pass away as soon as the child feels himself no longer under the sway of coercion. Far more effective than numerous prohibitions and commands, is the indirect influence of the parents’ attitudes upon the child: an exclamation of horror at moral vileness, unsimulated heartache and tears of sorrow over sin, sincere rejoicing in all that is good — these attitudes, repeatedly displayed, can not but produce in the child pleasure over right and displeasure over sin and thus prompt the will to action in the right direction. Only in exceptional cases will it be necessary to support such influence by threat or promise, punishment or reward. When the child has become conscious of the permanent presence of God to such an extent that his soul is filled with living confidence and reverence, then the most direct and unfailing way has been prepared for an initial recognition of sin, in fact for the quickening of all moral and religious feelings. But when these feelings have been stirred, an outlet for action must be provided; else they had better been left dormant. When, e.g., parental sympathy with the grief of others, or possibly the experience of tender help in his own needs, has kindled in the child, so self-centered by nature, a feeling of sympathy, then he must be shown a way of rendering help where actually needed and he must be encouraged to do it — even though it be no more than a flower that he offers to the afflicted person. If such influence is exerted upon him with regularity, then such
\quellseite{271}
actions on his part will become habitual and eventually will be the natural response to such situations. There is hardly anything that compares in importance with the formation of good habits in early childhood. It is true that this part of training, pertaining as it is to the pre-school period, must be carried out by the home. But this does not relieve the teacher of all responsibility; he will rather employ every means at his disposal to instruct the parents concerning their educational duties and unceasingly insist upon their performance.

\section{2. Childhood}

The entrance of the child upon the stage of boyhood and girlhood which takes place about the sixth year, is outwardly not so clearly marked as the transition from earlier (first year) to later infancy (second to sixth year), and appears rather like a continuation of infancy than something new. Yet unmistakable signs of a new period are not wanting. Physically, the change of teeth takes place at this time; muscles and nerves develop with greater rapidity; the brain has about reached its final size. Mentally, the child has made such progress that he can effectually participate in instruction. The fundamental period for the absorption of knowledge commences. Beginning to attend school he enters a new world by which he is unconsciously moulded and transformed, but which eventually he would, and should, also conquer. The period is likewise divided into two parts (from the sixth to the ninth and from the ninth to the twelfth year respectively) which we shall discuss separately.

a. Bodily, the restlessness of the child has not been overcome when, at six years, he begins to go to school: sitting still is so difficult that it should not be expected for any considerable length of time. But inasmuch as striving after activity no longer finds vent in aimless moving about from place to place as it did in the preceding period, but is naturally directed at definite objects, it is possible to arrange a curriculum
\quellseite{272}
that meets the requirements of this stage. Brief periods (no longer than 30 minutes) of actual impartation of knowledge should be followed by periods of technical instruction (writing, singing, drawing, manual work), and there should be numerous recess periods. Brief, frequently changing periods are required, particularly in the first year of school, by the mental development of the child, for his nervous and mental powers are still easily fatigued, and interest and attention can not long be focused on any one object. Sensation, perception, and \emph{intuition} are now the dominating factors of the intellect. The child of six lacks understanding for book-like accuracy almost altogether, and even the child of nine cannot appreciate abstractions. The desire to ask questions is rather heightened; all questions must be answered in concrete form and with constant attention to the peculiar concept world of the children, if an intelligent reaction is to be expected. The child wants to see, to behold. If he cannot see with the bodily eye—nine-tenths of our concepts are based upon sensory perceptions and intuitions—he desires to see with the mental eye whatever is put before his mind in intelligible, concrete form. Whoever, therefore, would successfully impart instruction in the period under consideration, must be a master in story-telling. And because the children usually have rather fragmentary concepts, even of the objects and conditions of their immediate environment; because their range of expressions is quite limited; and because they often attach a meaning to a word quite different from that of the teacher's; they will reap the full benefit of the teacher's talent of story-telling only if he first thoroughly investigates their verbal range and their sphere of concepts and if he raises their fragmentary and inaccurate concepts to the level of accuracy, clearness, and completeness. The faithfulness with which this duty is performed, will be reflected by the teachableness of the child, for the treasure
\quellseite{273}
of concepts thus acquired becomes the key with which the child unlocks the new material.

This capacity and this desire to behold things, however, are not the only factors demanding a method of instruction based upon plain, animated narration. This method is demanded also by the child’s \emph{phantasy} which is particularly active at this stage; not, it is true, the abstracting, but the constructing phantasy. With vigor and eagerness it stretches forth to grasp the concrete, to shape it according to its pleasure and to combine it into new forms as real as flesh and blood. The lifeless doll is transformed into a living, speaking, acting friend; the child lends it movement, growth, responsibility. Children would hardly find so much pleasure in play if it were not for the fact that it offers a field of activity not only for their impulse of action but also for their phantasy. Gradually the child’s phantasy emancipates itself from the limitations of time and space; it is no longer restricted to home and present, but transports itself into the remote past and into distant lands even though it may still array the persons living there with the concept material of home and present. But no matter whether phantasy transports itself, always the living person is its object of interest, and of imitation. As a matter of fact, the impulse of \emph{imitation} is at this stage of development very closely linked with the pronounced play of the child’s phantasy—another instance corroborating the demand to influence the soul of the child chiefly by narratives. Living persons, repeatedly presented to the phantasy, leave their traces in the soul; matters which ever and again have delighted the phantasy, are engraved upon the soul, especially during this “formative period of the brain.” The feelings being pleasurably excited, the will is aroused to imitation: the child would become similar to the person causing such pleasurable experiences.
\quellseite{274}
The \emph{emotions} as a whole, still center about self; unselfish friendships are rarely formed during these years, and if they are formed they may quickly be disrupted the moment the friend threatens to obscure the splendor of the own little Ego. Purposeful cultivation of the \emph{phantasy}, leading the child most clearly to visualize the grief of others, will seldom fail at this age to arouse the sympathetic \emph{emotions} and thus to put the will into action. As far as the \emph{volitional life} of the child at this stage is concerned, this period even more so than that of infancy, is of vast significance for the formation of \emph{habits}. Finally, it may be added that \emph{memory} though still weak at the age of six, from now on rapidly grows in strength; but it does not reach its greatest power in this period, and therefore only a limited amount of material should be assigned for verbatim memorizing, and all \emph{memory} material should at least be illustrated and, in a rudimentary way, be explained.

The conclusions to be drawn from these facts by the teacher of religion are self-evident. He must under all circumstances be thoroughly conversant with the concept world and the vocabulary of his pupils. Proceeding from the material really mastered by them, and if need be correcting it, he must aim slowly and gradually to enlarge and to enrich it. Biblical pictures will prove very helpful in this respect provided the teacher pays due attention to his diction, avoids all abstractions, forms only short sentences, ascertains by means of questions that he is really understood, engages the cooperation of his pupils, and trains them to express with their own words in simple sentences whatever he has taught them. Above all, he will, as soon as his pupils possess a sufficient vocabulary and a fair amount of clear and distinct concepts, employ narration, the most efficient means of instruction. His chief object must be to tell the story as simply as possible and at the same time as concretely and intuitionally as
\quellseite{275}
possible. On the one hand, he will most carefully bring out all the individual features, thereby securing life-likeness: this is what the child craves; and only thus are clear intuitions formed (cf. p. 204). On the other hand, he must stress the fundamental features of the particular story as a whole in order that the distinctive traits of the new story be not obliterated but stand out in clear relief even if to some extent the contents of the new story overlap with that of stories previously heard. Neither will he be satisfied with presenting merely the external features and their connection, but he will introduce his pupils into the hidden thoughts and motives of the acting persons in order to put into motion their awakening thought processes. He will confine his efforts to the most complete presentation of the main truth of any one story instead of developing from it a large number of less important truths, just as the unit of instruction at this time must be one simple story and not a group of stories fused into a unit. Extensive application does more harm than good at this age; but it is desirable to compress into some precious passage or a stanza the truth which has been discovered. The growing power of memory suggests not only the permissibility but the effectiveness of assigning such verses or stanzas as memory work (after repeating them in class with the children a number of times in unison); but the teacher will not foist upon the child's memory any material on which no proper light has been shed by preceding instruction. If it is true that no story, told well and repeatedly,\footnote{Stanley Hall says: “Of all the things that a teacher should know how to do, the most important, without any exception, is to know how to tell a story.”} ever remains entirely without influence upon the little child, but rather has the tendency and possibly to some extent also the power, to transform him into a likeness of its hero, how much more is this bound to be true of the Bible stories, especially of the stories of Jesus, and of the baptized child,
\quellseite{276}
cf. 2 Cor. 3:18. This effect is proportional to the measure in which the teacher succeeds in arousing the emotions of the child. The teacher may make the start with stories from present day life, but that the Bible stories may well be employed at this time, is evident from the play of phantasy which as was seen above transports itself readily into distant countries and remote times.—The teacher must also aim to stimulate the sympathetic emotions, and if he has influenced the will of the child, he must provide some opportunity for action; even at this stage, e.g., missionary zeal should be aroused and manifested.—In regard to discipline, the teacher will of course inculcate good order, so important for the later periods of life, but he must not forget that silence and sitting still are as yet difficult for the child. Against mendacity he will proceed with great earnestness, but because of the presence in these years of many incorrect concepts of the imagination, he must be warned that not every objective untruth is necessarily a lie and deserving to be treated as such.

b. The time between the ninth and twelfth or thirteenth year forms the second period of childhood. Bodily development usually progresses rather slowly during these years. Increase of stature and weight, compared with the rapid development in the preceding and following periods, takes place so slowly as to make it appear that progress, if not arrested, is at least seriously retarded. The growth of the brain, likewise, is quite inconsiderable. Only the connective tissue and, perhaps, the nerves exhibit vigorous growth. It seems as if nature is gathering strength for the sweeping changes of the subsequent period. Such accumulation of a store of energy and force is probably the cause of the noisy movements characteristic of boys of this age, and of their fondness for scuffling. For the same reason probably they prefer at this time physical to mental labor.
\quellseite{277}
As regards the \emph{mental life} of the children, the absorption of things and concepts approaching them from without still occupies the foreground. However, now the soul is not only capable of much more copious absorption, but also of a more thorough and accurate examination of the things absorbed—a step of progress which, in turn, results in a great improvement of the intuition. The organs of sense have not only become more delicate and discriminating, but the mental possessions acquired also prove an efficient key to the apprehension of the new and its coordination with the old (p. 212). The multitude of the impressions and concepts of which the soul at this time is capable, provided they are presented with concrete vividness, is astounding. Even more surprising, perhaps, is the ability to retain them. Between the eleventh and the twelfth year \emph{memory} develops its greatest strength and maintains its high level for several years afterward. This, of course, must not be understood to mean that the child at this stage learns more easily and rapidly than the adult, who has other and stronger aids at his disposal, especially the power of logical association; but the memory of the child is more faithful and accurate than that of the adult. His brain cells are still soft so that the impressions are deeper than those of the adult, whose brain cells permit only feeble impressions. For the literal memorizing of material of importance for the whole life, the eleventh and the twelfth years represent “the golden period.” Foreign languages if they are really to be mastered, are most readily and accurately acquired down to the end of the fourteenth year. The achievements of memory increase in direct proportion to the growth of attention, interest, and reasoning, and to the accumulation of apperceptive material. Experience and experiments in regard to the powers of memory have shown that it is more profitable to rest a while between repetitions than to repeat the subject; say, eight times in quick succession; and, likewise
\quellseite{278}
to pause after the study of one subject than at once to proceed to the study of another. The inner processes connected with learning continue to act for some time after the conscious processes have ceased. The teacher must take account of this retroactive inhibition: if fresh impressions disturb these mental processes, memory is bound to suffer for the reason that the new impressions interfere with what has just been learned, and especially so when close attention is paid to these new impressions.

The \emph{phantasy} continues very active at the age under consideration, requiring food and cultivation; but as experience and judgment attain to greater maturity, fondness for the impossible is stripped off more and more. Accordingly it is the exception when the fairy tale still fascinates the soul. But wherever realities are placed before the soul, particularly in geography and history, the phantasy, aided by charts and pictures, builds for itself a rather correct image of distant scenes and times, especially if it be stirred by vivid descriptions. Also at this stage of development the child requires concrete descriptions; but these should accurately agree with reality, the time, place, and sequence of events. While the phantasy and the sense of history stand correlated, it is the personal, biographical element in history that most attracts the soul. No interest is shown as yet in institutions, in forms of organization and in laws except through vital connection with great men and women. The history of a whole nation, likewise, is imparted to children of this age only through the presentation of its outstanding heroes. This is the age of hero worship. The stronger and bolder, the more clever and daring these heroes prove to be, the more assuredly they will become the objects of devotion, especially for the boys. Wherever there is fondness for reading at this age, preference is given to hero stories both good and bad.—The \emph{thinking activity} is marked by decided progress during these
\quellseite{279}
years. In the same measure in which the child becomes able to observe the objects of the external world with greater accuracy, and to distinguish the essential from the unessential and accidental, he learns to coordinate the essential marks and to form conceptions, to affirm judgments, and to draw conclusions. Already, though hesitatingly, questions in regard to the connection of things are asked. Such questions, at first, refer only to objects of the nearest environment, for instance a toy or watch which the boy takes apart in order to acquaint himself with this marvelous mechanism, and to natural phenomena which lie before his eyes like baffling riddles; soon, however, the mere sequence of events both in human life and in history ceases to satisfy him. The reason of the child begins to inquire in regard to cause and effect. What is wanted is not necessarily a scientific proof; a reasonable explanation, if it is only plausible, is accepted upon authority. Not seldom such initial critical activity of the young mind invades the religious sphere.—From what has been said, it is evident that there is a strengthening in these years of the \emph{intellectual feelings}. The esthetic feelings, too, under proper guidance, awake at this time. The ethical emotions assume more definite forms; wherever the development has been normal, conscience now enters as a strong factor although suggestion and the example of adults are as yet still more powerful. God, especially by boys, is felt as Law and Power; but He can be experienced also as Holy Love, especially by the baptized child, so that he does not only view God with awe but surrenders to Him with all the power of his youthful emotions. The social feelings are now quickened as never before. Hence this is the period for the formation of friendships: the child associates with others of the same age, disposition, and condition in life. “The bonfire with its circle of kindred spirits, the cave with its password and dark plottings, the street corner
\quellseite{280}
and recruiting whistle, have almost irresistible fascination.” Even the girls have their own clubs. While truth and loyalty are observed toward those of the inner circle, the child, at this age, often follows a peculiar code of morals toward those outside. Self-esteem and the sense of honor are conspicuous in their development. The emotions as a whole still change often and swiftly; remorse or sadness does not last long at this age: happiness is the dominant note.

— Of the greatest significance for this period, finally, is the formation of good \emph{habits}. While the conditions therefore have existed in the previous period, they are now most favorable, more so than they will ever be again. Thoughts that are repeatedly entertained, frequently recurrent emotions, decisions of will formed time and again, the moral acts thus resulting—all these leave in the soul impressions or grooves, as it were, along which the subsequent activities of the soul run almost automatically; and it becomes increasingly difficult for the child to think, feel, and will in any other direction. Therefore, if the good, the noble, the true, if God and His will, at this time, determine the direction of the soul, the result will be that, in normal circumstances, no special reflection, no special conflict is needed before the good act is performed: to act in a given way has become the second nature of the soul.

This is also the time when the \emph{difference of sex} becomes more clearly established. The girl usually has a more delicate constitution than the boy, both of body and soul, requiring for this reason a more delicate treatment. The boy is more awkward of movement, more abrupt, impetuous, inclined toward physical exercise and boisterous play; the girl has more nimbleness, gracefulness, harmony of movement, and is gentler and more retiring. This difference is plainly evident in the games played by boys or girls. Mentally the girl, as a rule, is swifter and more agile, but less thorough and
\quellseite{281}
keen than the boy. The preponderant emotions of the boy are the sense of fairness and self-esteem, and interest in things intellectual, great, sublime, heroic; of the girl, sympathy, the sense for the beautiful, lovely, proper. Similarly, in point of good recitation, diction, and handwriting, the girl is superior to the boy. The boy, above all, expects from his teacher justice while the girl looks for kindly treatment. The will of the boy excels that of the girl in resoluteness and force, but is inclined to stubbornness, impetuousness, and lack of consideration; the girl is more tender-hearted, patient, and tractable than the boy, inclining, however, more than her brother, to sensitiveness and vanity.—These years also witness the establishment of the pupils’ \emph{individuality}. In general aspects quite similar to his companions, each pupil possesses distinct and individual characteristics, intellectually, emotionally, and volitionally. To study and to consider his pupils’ individualities according to their powers and weaknesses is an indispensable duty of the teacher.

All these facts must be carefully borne in mind by the teacher of religion. What has been said concerning bodily development will prevent him from adjudging as wilful opposition to good order and discipline when his pupils indulge in noise and commotion. While he will endeavor to keep the desire for scuffling in bounds, he will not suppress it for fear of choking the germs of promptings that, later, might prove of eminent value. “If you crush the fighting instinct, you get the coward; if you let it grow wild, you have the bully; if you train it, you have the strong, self-controlled man of will.” The great capacity of the pupil for copiously absorbing impressions and concepts, and to retain them by the power of his memory, suggests to the teacher the right and the duty to utilize these years as the great \emph{period of learning} also from the religious aspect. He will let his pupils memorize Scripture passages, hymn stanzas, and sections of
\quellseite{282}
the Catechism; but never, be it understood, without placing the memory material into the proper light or at least explaining its verbal meaning. The proper elucidation of memory material is particularly valuable since, irrespective of other advantages, it arouses interest in the material to be memorized and employs association as an aid to appropriation. He will also show the child how to memorize: how to analyze the whole into sections as suggested by the contents, and how to observe pauses in the learning process in order to give the impression an opportunity to settle before the brain is assailed anew. He will make the \emph{phantasy} ancillary to instruction by making the pupil visualize concrete and plastic forms from the history of the Bible, of the Church, and of missions, and by adjusting himself to the pupil’s craving for reality through the diligent use of biblical and other charts and every sort of available pictorial representation. He will not deem it a loss of time to dwell for a rather long time upon some heroic personages and to pursue their life story from birth to death. The short narrative, at this stage, is not imperative any more. The better the child can pursue the development of the hero as a whole, especially the years of his youth, the more certainly he will actually live with the hero, enter into his experiences, feel sympathy with him. Thus sets in the eminently important “ideal association” with the hero whose life is studied. The development of the \emph{thinking faculty} now permits or, rather, prompts the teacher to enter with his pupil into the inner motives, the impelling powers of such a life, to exhibit the relation between disposition and action, and to set forth the laws of the moral and religious life; in this way he will train him to form correct moral and religious judgments and arouse the desire and the courage to apply them to his own conduct. This is the teacher’s plain duty; he must make the application to the life of the pupil and thus exert a \emph{moulding influence} upon
\quellseite{283}
it although the application made by him or, at his direction, by the pupil, should never be dragged in but rather naturally grow forth from the subject matter and its treatment by the teacher; it should impress itself upon the pupil by inner necessity. When the teacher has really touched the pupil’s feelings and has brought about a decision of will, he must make the further effort to bring about the performance of the deed and the establishment of \emph{habit}. The life at home, at school, in the congregation, among the neighbors, should be set forth as the sphere of moral and religious activity. Veneration for age, respect for the powers that be, obedience, conciliatoriness and self-control, sympathy with others, helpfulness, self-denial, devotion, sacrifice, consecration before the eyes of the ominscient and omnipresent One, prayer to the great heavenly Friend, all these attitudes must become matters of \emph{habit}, now more than ever they were before. The practise of church attendance must now begin. Salutary even before this age, there are particular reasons why it should now become habitual. While the children are not likely to understand the sermon as a whole, they will understand sections of it. To this factor must be added the sight of the congregation at worship, the solemn liturgy, the power of song, and the impressions gained from the pictures that may belong to the equipment of the church—all this cannot but permanently affect the soul. The teacher will appeal to the thinking faculty of the pupil also for the purpose of giving him tasks that compel thought; of training him to form important religious and moral conceptions, conclusions, and judgments, always be it remembered, upon the basis of the intuition inasmuch as the pupil otherwise would have no access to conceptions and abstractions. Let the teacher of religion show a sympathetic attitude toward the critical questions of a religious character which, at this stage, will occasionally be propounded to him. Such questions, as a
\quellseite{284}
rule, arise from a rather anthropomorphical view of God and of the invisible realities in general. Still better, let him anticipate them by taking the initiative in stripping off such anthropomorphic features from the pupil’s views of God and heaven, making it plain to him that God is a Spirit; that heaven is complete union with God; and that things corporeal and, therefore, the objects of space and time are only pictures and symbols of invisible realities of which men cannot speak save in material terms. If, however, such critical questions show doubt as to the existence itself of these heavenly realities, then the teacher may remind the young critic of the as yet very limited range of his knowledge and experience, or better, he will by his own conduct, especially his happiness in fellowship with God, most effectively remove such doubts. It may be stated as a general truth that living models and examples are bound to show themselves as exceedingly effective educational factors. For this reason the teacher will keep an eye upon the friendships bound to be formed at this stage. He will exhort the home to exercise vigilance in this respect and to contribute as much as possible to the gratification of these social feelings by sweetening family life and by occasional invitations of congenial children. — Finally, when the child is soon restored to happiness after bereavement has entered the home or after a lapse into sin, the teacher will not conclude that the child is frivolous or unfeeling; what he observes is in keeping with the child’s nature. And although it is surely true that the Holy Spirit can lead a child to true repentance of sin and an unqualified surrender to Jesus Christ at this stage of his life, the teacher of religion will take care not to apply the same criterion to these processes in the child soul which he is wont to apply to them in the soul of the adult. He would only be guilty of injustice or lead his children to perilous hypocrisy.

\quellseite{285}
\section{3. The age of adolescence}

With the thirteenth year, the girl; with the fourteenth, the boy, enters upon the age of youth. In the case of the girl, this period, extending to the 20. or 21. year of life, comprises the whole process of development from girlhood to womanhood; in the case of the male, this period, extending down to the 23. or 24. year, involves the whole time of development from boyhood to manhood. These bare facts indicate that this is the period of the most revolutionary change. Stanley Hall, the pioneer among its observers, does not overestimate its significance when he boldly calls it a “new birth.” There is a new birth of the body and of the mind, the one most intimately linked with the other.

More fundamental than ever is the \emph{bodily change} now taking place. The body, usually growing with rapidity, increases in height and breadth; the bones become stronger and are joined more firmly; the limbs lengthen; especially the hands and feet assume dimensions out of all proportion to the rest of the body; the heart increases its size one fifth and, in some cases, even more; blood pressure rises materially; respiration becomes one-third slower and, in consequence, more intensive. While the brain has in general attained its final size, certain centers in the cerebral membrane (meninx) do not begin to function until now; the tissue involved becomes denser and promotes a slow growth of the cranium. The most important change, however, is the attainment of \emph{sexual maturity}. The sexual organs assume their final form; in the male there begins the occasional emission of semen, in the female, menstruation, and, in connection therewith, the development of the breast; in both, the appearance of hair at places hitherto bare, and the change of voice. The change is more conspicuous in the female and is often accompanied by violent pain; but also in the male it often proceeds amid far-reaching constitutional disturbances that may affect the whole organism.
\quellseite{286}
Both sexes often suffer from anaemia, because the development of lungs, liver, kidneys, and the muscular system, the entering of the sexual organs upon their functions, and an increased activity of the central nervous system require a greater supply of blood. The excitation of the organism may, in the case of both, go to the length of producing a suspension of consciousness, temporary dementia, or even a complete shattering of the nerves and ultimate insanity. But even where this change takes place with comparative smoothness, it brings about a situation altogether new.

Hand in hand with this physical change goes the \emph{mental}, which is largely based upon the other. The vigorous development of the muscles and the whole body calls forth, especially in the maturing male, a sensation of \emph{strength}, as soon as the limbs have gained their ultimate proportions and he has ceased to fall over his own feet. Pityingly he looks upon the childhood stage which he has hardly left behind, and fancies that he is already a man who has reached the height of his power and knows of no obstacle that he is not able to conquer. This feeling of strength and the urge of his restless blood account for his clownishness and turn him into a daredevil for whom no venture requiring strength and courage is too great. Having become a man, as he thinks, he means to be \emph{independent}, to be free from all dictation save his own, at least to raise a mustache, to smoke, perhaps to gamble and to curse just like a man. Whoever tries to prevent the youth from indulging his whims, is regarded as the foe of all true progress. The passion for liberty may assume a revolutionary form and array itself in opposition to every authority and barrier. If the social conditions are such that the male youth can earn money for himself, the feeling of independence often assumes the most repugnant forms and easily leads to a break with the parental home. But when the young man is made to experience the limits of
\quellseite{287}
his strength and power, he may not only seek to replenish them by excessive eating and sleeping, but his self-reliance and independence may suddenly change into sheer \emph{despair and dejection}—a condition in which he needs nothing so much as a real man or an understanding mother heart that will restore his courage.—Quite similar is the development of the female. To the boy clown corresponds the girl clown with her perpetual giggling and saucy conversation. Instead of muscular exercises, she prefers games requiring agility and speed of movement. In point of self-reliance the girl of this age is likely in no wise to remain behind her brother: she wants to dress and to deport herself like a mature woman and to be treated accordingly.

Contemporaneous with the change of the cerebral membrane and, perhaps, in some measure occasioned by it, there is a prevailing inclination during these years toward logical thinking. This fact explains to a certain extent why youth, despite a lack of real knowledge and adequate experience, claims to possess a knowledge of everything, and usually a \emph{knowledge superior to that of adults}, and that he deems himself warranted, or even summoned, to criticize his parents and teachers. Other factors contributing to this condition are possibly the aimless and omnivorous consumption of all sorts of questionable or, at any rate, insipid literature, or the fact that the school has been outgrown, or that the teachers in school did not insist on thorough comprehension of the problems under consideration (esp. those of the natural sciences). But this is only one phase of the adolescent mind. Often the young deeply feel their immaturity, the inadequacy of their knowledge, and by no means only when they have been made aware of it by some painful experience; sometimes they feel it most keenly when they attempt to disguise it through a bombastic attitude, loudness of voice, the unintelligent use of axioms, appeal to authorities—with all their
\quellseite{288}
vaunted independence they are fond of quoting authorities—, and if need be, the use of the fist. What the adolescents need is a man of thorough knowledge and loving sympathy who will recognize their eagerness to learn, put to work their thinking apparatus, and lead them to mature knowledge. To him they will be devoted especially if he imparts a type of knowledge which possesses practical importance for the daily life and pertains to their environment. This applies with double force to the rural youth, so often devoid of mental initiative, rather slow in formulating independent thoughts and questions, and withal very bashful in the presence of pastor or teacher.

The most important feature of the inner life of the young, however, stands connected with the \emph{development of puberty}. All at once the soul becomes aware of new, mysterious, hitherto unknown forces and impulses in the organism. Exercising powerful attraction, the element of mystery brings the soul under its spell. The impulse to penetrate the mystery and to learn something reliable concerning its origin and functions is awake. The soul begins to presage the truth; it begins to understand itself from the aspect of sex; the relation to the other sex undergoes a change; the former ease in the presence of the other sex has disappeared; the \emph{attraction of the one for other} makes itself felt—sooner with the girl, however, than with boy, who, if morally healthy, will feel, and readily exhibit, toward the girl something very much akin to contempt. At times, genuine youthful love awakens, here and there also enthusiastic affection for persons far older than oneself. The soul begins to occupy itself with the future which apparently teems with possibilities. The \emph{phantasy}, probably in connection with the above mentioned changes in the brain, unfolds a markedly vigorous activity and paints the future in glowing colors. The soul becomes a dreamer. There is a heightened understanding
\quellseite{289}
for poetry. The desire to read is in evidence; nothing attracts the girl more than the love story while the boy yields to its fascination less readily except when an all-conquering hero is in the center of the plot. What appeals to him is not so much sentiment as courage, daring, heroism. When the soul, provided that it is uncorrupt, comes in contact with the good, noble, great, it reaches out after it with longing; \emph{ideals} are formed which it means to realize never doubting that such achievement is within the range of its knowledge and power. Nor is the soul any longer self-centered: the \emph{social feelings} begin to exert their power; the desire is felt to help, correct, reform, the world. But if during this period the soul falls into the hands of wicked friends or, if it feeds upon salacious literature, it may be set on fire in behalf of the questionable, low, destructive. Often the two lie side by side; youth is a child of the moment, prone to harbor contradictions. After the manner of spring, it makes progress by contrasts, not necessarily religious and moral, but temperamental. Now heavenly joy, then deathlike grief; now excess of enthusiasm, then of despondency to the point of despair; now the delight over hearth and home, then the yearning for the everlasting snow of lofty mountains and \emph{Weltschmerz}, grief for the world's grief; now light-hearted candor and frankness, then somber reserve and desire for solitude—such contrasts as these, largely in response to physical states, often supersede each other in rapid sequence. Still another factor merits consideration in this connection: the sexual impulse by this time is fully aroused and, in the state of depravity, threatens to assume the role of a tyrant; but the soul is aware or has at least a presentiment of the fact that there dare be no yielding, not even in the form of masturbation. Thus a \emph{conflict} is occasioned between soul and body which pervades the whole human organism. While this battle is in progress, one phenomenon in particular is in evidence:
\quellseite{290}
conceptions such as conscience, God and the divine sanctions, sanctification, responsibility, likewise sin and grace receive new or deeper meaning. The soul strives for nothing so much as for moral strength; where that is not found, or where the soul has been so bestialized that it never even enters upon the conflict, it may yield in helpless abandon to the exactions of desire, presently to stand at the brink of despair or to go down in it altogether. From this point of view, and from others previously mentioned, it is not surprising that moral self-decision, genuine conversion, or a deepening of the inner life is so frequently experienced in connection with the transition from adolescence to puberty (chiefly in the sixteenth and seventeenth years); just as the beginning of the downward trend may often be traced to this period. We are face to face here with a crisis of the most transcendent importance from which man comes forth bodily, mentally, and sometimes even spiritually, as one born anew so that often the parents themselves marvel at the change undergone by their child, — a crisis in which also the Spirit of God is often at work upon the soul. — The difference between the young man and the young woman is found principally in this that in the latter the emotions rather than reason and will develop into the dominant force, a difference which is bound to remain despite some foolish tendencies in the “woman movement” which appear to set in the opposite direction.

Youth being thus constituted, whoever would become its leader in the moral and religious sphere must be ready at the outset to leave the rights of youth uncircumscribed, and to take due account of them in all educational measures. Whoever would treat adolescents as if they were devoid of power or knowledge or ability, in short, as children, will see them turn away from him with frankest repugnance. Whoever would keep them in bounds by a formal appeal to authority may experience a revolt at their hands. Whoever
\quellseite{291}
lacks acquaintance with their problems and needs or disdains to occupy himself with these will find their hearts closed against himself from the beginning. Whoever would train them as hermits or monks will never behold their eyes lighting up with responsive gleams; for joy and mirth is youth’s heritage. Whoever would stamp as sin what is not sin will find unwilling ears even when he rightly takes exception and attempts to invigorate their conscience. Whoever indulges in dogmatizing will cause the young ennui; and ennui is youth’s mortal foe. Whoever, finally, has flattened religion to a matter of mere forms will see the young turn their backs upon him because they are hungry for life, or because their innate sense of truth revolts against such a \emph{formal religion} and condemns it as untruth.— But he who acknowledges the right of the young to progress, and stands by their side in times of helplessness and despondence as a trusting friend, will win their confidence and they will cling to him as the vine clings to the pillar. Whoever does not look down with contempt upon such knowledge as the young possess, but gratifies their thinking faculty by the imposition of worthy tasks, thus supplying it with incentives to effort; whoever is not too proud to learn with the young and does not try to cover an occasional gap in his own information with empty words or, still worse, with pious phrases; whoever meets their misgivings and doubts with sympathy and then, step by step, conquers them with real arguments, to such a teacher they will be devoted. Whoever impresses the young that he is religious, not because he has remained unacquainted with science, with the leaders in the world of thought, with literature and modern problems; but having acquired knowledge, thorough knowledge, of them, having digested them, has remained religious notwithstanding, he is their man: to him they will open ear and heart. He who in reality is a Christian character, a well-spring of power and life, open to the
\quellseite{292}
world and yet fettered to Christ, he alone can be a guide to youth, which soars high one moment only to sink into the slough of despond the next. A protector of youth’s freedom, he yet constrains them to follow him. Sure to disclose Christ to them as Saviour, he will yet stress the fact that this Saviour is a hero of strength to whom one can look up as an example and from whom power proceeds to conform to the model. Such a teacher will meet the social impulse of the young and organize them in a society, remembering, however, that there is a gap in their development, in the sixteenth year in the case of the girls, in the seventeenth or eighteenth in the case of the boys, and that from these years onward they must be allowed more independence and freedom. Moreover, he will remind himself that any organization will collapse from inner weakness unless opportunity, coupled with leadership, be afforded for action.
